% !TeX root = ../../Book4.tex
% 纲要标题：## 第5章　势论、统一化与双曲曲面
% 纲要位置：几何学计划纲要.md，第 3629 行。
% 正文状态：初稿；尚未正式审核。

% 纲要内容开始
% 
% > 本章目标与证明范围：为一般统一化保留独立而连续的证明模块：先做带边界曲面上的势论，再用拓扑穷竭及正规族处理单连通非紧曲面。不能只引用平面映照定理后宣称所有曲面均已统一化。
% 
% 纲要内容结束

\chapter{势论、统一化与双曲曲面}
\label{g4:ch:05}
\BookChapterNumber{g4:ch:05}

\begin{chapterabstract}
	局部复坐标如何拼成整张曲面的标准模型，是统一化要解决的问题。本章从带边界区域的调和函数开始，以 Poisson 替换和边界障碍证明 Dirichlet 问题可解，再构造具有对数奇点的 Green 函数。对单连通带边界区域，Green 函数的负指数给出固有的一次圆盘映照。一般非紧单连通曲面则先经填洞得到光滑圆盘穷竭，再用归一化导数、单叶函数的畸变估计及逆映射紧性得到整个曲面的映照。由此证明球面、平面和圆盘的统一化分类，并从万有覆盖得到一般曲面的商描述。最后构造双曲型曲面的完备 Poincaré 度量，证明负曲率下的收缩，计算尖点、漏斗、极值长度与长颈模量。
\end{chapterabstract}

\begin{learninggoals}
	\begin{itemize}
		\item 用局部 Poisson 积分、次调和函数的上包络和边界障碍，完整求解光滑相对紧区域的 Dirichlet 问题。
		\item 构造 Green 函数，计算共轭的周期，并以固有性和映射次数证明圆盘映照的双全纯性。
		\item 在非紧单连通曲面中证明光滑圆盘穷竭，把归一化映照的尺度与共形半径联系起来。
		\item 推导单叶函数的增长与四分之一圆盘估计，分别控制映照及其逆，证明统一化极限非退化且满射。
		\item 计算曲率的共形变换，证明 Ahlfors--Schwarz 比较与 Poincaré 度量的完备性、收缩及等号条件。
		\item 计算环域与曲线族的极值长度，用模量、核心长度和截口参数描述具体颈部的退化。
	\end{itemize}
\end{learninggoals}

平面 \(\mathbb C\) 与圆盘 \(\mathbb D\) 都是单连通曲面，局部坐标也没有区别，它们却不可能双全纯同构：任何从整个平面到圆盘的全纯函数都有界，因而为常数。局部复结构相同并不决定全局共形类型。另一方面，穿孔平面 \(\mathbb C^\times\) 与穿孔圆盘有相同的无限循环基本群，它们的万有覆盖却分别是平面与圆盘。要辨认曲面的类型，需要同时控制函数的全局延拓和它在远处的尺度。

\ProseParagraphBreak

先把问题限制在一块有边界的曲面区域中。单位圆盘上的函数 \(-\log|z|\) 在边界为零，在中心有一个对数奇点；局部共轭使它成为 \(-\log z\) 的实部，取负指数便恢复坐标 \(z\)。如果在抽象曲面区域上也能造出具有这些性质的函数，同样的指数构造就能给出圆盘坐标。单连通性控制绕路造成的整数周期，边界条件控制原像的紧性，一个简单零点再把映射次数固定为 \(1\)。势论因此提供了从实函数进入全纯坐标的途径。

\ProseParagraphBreak

对没有边界的非紧曲面，需要让这些有限区域逐渐扩大。每一步的映照都可以把同一个基点送到原点，但若只作这一点归一化，极限仍可能变成常数。我们再把该点的导数固定为 \(1\)，让目标圆盘的半径记录尺度。当半径趋于有限值时，极限目标是圆盘；当它趋于无穷时，目标是整个平面。证明必须同时保证正映射有局部一致界、逆映射不会逃出曲面的紧集，这正是单叶函数的畸变控制与四分之一圆盘估计承担的任务。

\ProseParagraphBreak

本章使用第一章的局部全纯函数理论和正规族、第三章的亏格零情形，以及第一卷的紧曲面分类、覆盖提升和光滑单位分解。所需的拓扑圆盘穷竭、势论和单叶函数估计都在本章展开。平面上的 Harnack 原理与 Perron 方法可对照\parencite[第6章第3.2节、第4.2节]{Ahlfors1979Complex}；这里把局部计算组织在曲面上，再处理一般单连通曲面的穷竭。全文沿用非负 Laplace 算子 \(\Delta=-\operatorname{div}\nabla\)，双曲度量的曲率归一化为 \(-1\)。

% !TeX root = ../../Book4.tex
% 纲要标题：### 5.1 曲面上的势论工具
% 纲要位置：几何学计划纲要.md，第 3633 行。
% 正文状态：初稿；尚未正式审核。

% 纲要内容开始
% 
% * 辅助共形度量和调和性
% * 光滑相对紧区域的 Dirichlet 问题
% * Green 函数及对数奇点
% * 调和共轭、周期和能量
% * 最大值原理与 Harnack 紧性
% 
% 纲要内容结束

\section{曲面上的势论工具}
\label{g4:sec:0501}

要把一张曲面映到圆盘，可以先寻找一个实函数：它在边界取零，在一个指定内点像 \(-\log|z|\) 那样趋于正无穷。如果这个函数在其余地方调和，配上局部调和共轭再取指数，便有机会得到所需全纯映射。本节先建立这个实函数的存在性。所有局部计算都在全纯坐标中进行；度量用于组织能量与法向量，调和性本身只依赖复结构。

\subsection{辅助共形度量和调和性}
\label{g4:subsec:0501:01}

设 \(X\) 为黎曼曲面。全纯坐标变换 \(w=w(z)\) 满足 \(|\mathrm{d}w|^2=|w'(z)|^2|\mathrm{d}z|^2\)，所以这些局部 Euclid 度量彼此共形。第一卷第12.3节定理12.3.5给出从属于坐标覆盖的光滑单位分解：支集族局部有限、各支集留在对应图内、权重非负且总和为 \(1\)。把各图中的 \(|\mathrm{d}z|^2\) 乘权重后补零并求和，便得到全局光滑正定度量，在每张全纯图中形如
\[
	g=\rho(z)^2(\mathrm{d}x^2+\mathrm{d}y^2),\qquad \rho>0.
\]
局部有限性保证总和光滑，正性来自至少一个正权重。这称为与复结构相容的\term[gongxingduliang]{共形度量}[conformal metric]。

\begin{proposition}
	\label{g4:prop:0501:conformal-harmonicity}
	设 \(X\) 为黎曼曲面，\(g=\rho^2|\mathrm{d}z|^2\) 为光滑共形度量。按非负 Laplace 算子的约定，对实 \(C^2\) 函数有
	\begin{equation}
		\label{g4:eq:0501:laplacian}
		\Delta_g u=-\rho^{-2}(u_{xx}+u_{yy}).
	\end{equation}
	因此条件 \(\Delta_g u=0\) 与所选共形度量无关，并在全纯换坐标下不变。这样的函数局部为全纯函数的实部，且光滑。
\end{proposition}
\begin{proof}
	度量矩阵为 \(\rho^2I\)，面积密度为 \(\rho^2\)。散度的坐标表达给出
	\[
		-\operatorname{div}_g\nabla_g u
		=-\frac1{\rho^2}\sum_{j=1}^2
		\partial_j\bigl(\rho^2\rho^{-2}\partial_j u\bigr)
		=-\rho^{-2}(u_{xx}+u_{yy}).
	\]
	若 \(w=w(z)\) 全纯且导数非零，链式法则及 Cauchy--Riemann 方程给出
	\[
		(\partial_x^2+\partial_y^2)(v\circ w)
		=|w'(z)|^2(v_{ss}+v_{tt})\circ w,
		\qquad w=s+\mathrm{i}t.
	\]
	因而零解在坐标间相容。在每个小坐标矩形中，\cref{g4:prop:0103:harmonic-conjugate}沿两条坐标线的积分构造出共轭函数；实部来自全纯函数，故由局部幂级数具有任意阶导数。
\end{proof}

曲面的复结构还在实一形式上给出旋转算子
\[
	*\mathrm{d}x=\mathrm{d}y,\qquad *\mathrm{d}y=-\mathrm{d}x.
\]
全纯坐标变换的实微分是正伸缩与旋转的复合，所以它与此算子交换；\(*\) 在不同图中定义同一个全局算子，也不受共形因子 \(\rho\) 影响。于是
\begin{equation}
	\label{g4:eq:0501:star-laplacian}
	*\mathrm{d}u=-u_y\,\mathrm{d}x+u_x\,\mathrm{d}y,
	\qquad \mathrm{d}(*\mathrm{d}u)=-(\Delta_g u)\,\mathrm{d}A_g.
\end{equation}
调和函数的共轭问题因此成为闭一形式 \(*\mathrm{d}u\) 是否恰当的问题。以下只用调和方程的零解；需要判断次调和性时，采用的局部不等式是 \(u_{xx}+u_{yy}\ge0\)，即本书记号中的 \(\Delta_g u\le0\)。

\subsection{光滑相对紧区域的 Dirichlet 问题}
\label{g4:subsec:0501:02}

称 \(\Omega\Subset X\) 为光滑相对紧区域，是指 \(\Omega\) 开且连通，\(\overline\Omega\) 为嵌入 \(X\) 的紧光滑带边界曲面，并且 \(\partial\Omega\ne\varnothing\)。边界在每张小坐标图中是一条光滑曲线，\(\Omega\) 位于其一侧。\term[dirichletwenti]{Dirichlet 问题}[Dirichlet problem]\footnote{Johann Peter Gustav Lejeune Dirichlet（1805--1859），德国数学家，研究数论、Fourier 级数与势论。}要求对给定连续边值 \(f\in C(\partial\Omega,\mathbb R)\)，构造内部调和、连续到闭包且边值为 \(f\) 的函数。唯一性来自最大值原理，存在性则由下面的圆盘求解与上包络构造完成。圆盘上的显式解称为 Poisson 积分。\footnote{Siméon Denis Poisson（1781--1840），法国数学家，研究积分、级数与数学物理。}

\begin{lemma}
	\label{g4:lem:0501:poisson}
	设 \(f\) 为圆周 \(|z|=R\) 上的连续实函数。则
	\begin{equation}
		\label{g4:eq:0501:poisson}
		P[f](z)=\frac1{2\pi}\int_0^{2\pi}
		\frac{R^2-|z|^2}{|Re^{\mathrm{i}t}-z|^2}
		f(Re^{\mathrm{i}t})\,\mathrm{d}t,
		\qquad |z|<R,
	\end{equation}
	调和并连续延拓到闭圆盘，其边值为 \(f\)，且这样的延拓唯一。若 \(h\ge0\) 在 \(D(0,R)\) 上调和，则对 \(|z|\le r<R\) 有
	\begin{equation}
		\label{g4:eq:0501:harnack-disk}
		\frac{R-r}{R+r}h(0)\le h(z)
		\le\frac{R+r}{R-r}h(0).
	\end{equation}
	在每个较小闭圆盘上，调和函数的各阶导数由稍大圆盘上的函数一致界控制。
\end{lemma}
\begin{proof}
	核是全纯函数 \((Re^{\mathrm{i}t}+z)/(Re^{\mathrm{i}t}-z)\) 的实部，所以在紧子盘上可逐项求导，所得积分调和。核非负，且对 \(t\) 的平均为 \(1\)：将上述分式展开为 \(1+2\sum_{n\ge1}(z/R)^ne^{-\mathrm{i}nt}\)，非零频率积分为零。故 \(|P[f]|\le\max|f|\)。

	固定边界点 \(\zeta\)，把积分分为靠近 \(\zeta\) 的小弧及其补集。在小弧上，\(f-f(\zeta)\) 由一致连续性任意小；在补集上，\(z\to\zeta\) 时分母一致离零，分子趋零，所以该部分积分趋零。核的总积分为 \(2\pi\)，因而 \(P[f](z)\to f(\zeta)\)。唯一性是\cref{g4:cor:0103:harmonic-maximum}的应用。

	对 \(h\) 先在任一半径 \(R'<R\) 的闭盘上用唯一性得到 Poisson 表达。核在 \(|z|\le r\) 上介于 \((R'-r)/(R'+r)\) 与其倒数之间，而边值的平均为 \(h(0)\)。令 \(R'\uparrow R\)，便得式\eqref{g4:eq:0501:harnack-disk}。对 Poisson 核求任意阶实偏导，在固定内盘和固定较大积分圆周之间，分母有统一正下界，遂得到相应的导数估计。
\end{proof}

为作上包络，我们暂时允许尚未调和的连续函数。称连续实函数 \(v\) 为\term[citiaohehanshu]{次调和函数}[subharmonic function]，是指每个闭包位于定义域内的坐标圆盘 \(B\) 及每个 \(H\in C(\overline B)\)、内部调和的函数都满足：若 \(v\le H\) 于 \(\partial B\)，则 \(v\le H\) 于 \(B\)。这等价于每个小坐标圆周上的次平均不等式。一个方向用 \(P[v|_{\partial B}]\)；另一个方向注意，若 \(v-H\) 在盘内有正最大值，次平均不等式迫使最大值点附近恒定，再由连通性把最大值传到边界，产生矛盾。因此定义可由局部次平均性质检验。相同的最大值论证使比较性质对任意相对紧的小区域成立；全纯换坐标保持调和函数，故定义也与选图无关。\(-v\) 次调和时，称 \(v\) 为超调和函数。

\begin{lemma}
	\label{g4:lem:0501:perron-operations}
	连续次调和函数的和及有限最大值仍次调和。在一个坐标圆盘 \(B\Subset\Omega\) 内，把连续次调和函数 \(v\) 换成其边值的 Poisson 延拓，在盘外保留 \(v\)，所得函数记为 \(v^B\)。则 \(v^B\) 连续、次调和且 \(v^B\ge v\)。连续次调和函数在连通区域中若取得最大值，就恒为该常数；特别地，闭包紧且边界非空时，其上界由边界控制。
\end{lemma}
\begin{proof}
	次平均不等式可相加，故和仍次调和；有限最大值的比较性质逐项成立。次平均不等式也证明强最大值原则：在最大值点，圆周平均不小于中心值而又不大于全域最大值，故整圆周都等于最大值；让半径变化，最大值集合为开集，它又是闭集，连通性给出结论。对于只有局部最大值的情况，先在小盘中得到局部常值；全域边界比较时只需使用全域最大值。

	Poisson 延拓连续至边界，故 \(v^B\) 连续；次调和比较给出 \(v^B\ge v\)。检验接缝处的局部次平均性质，设 \(p\in\partial B\)。有 \(v^B(p)=v(p)\) 且邻域内 \(v^B\ge v\)，所以
	\[
		v^B(p)=v(p)\le\frac1{2\pi}\int_0^{2\pi}v(p+re^{\mathrm{i}t})\,\mathrm{d}t
		\le\frac1{2\pi}\int_0^{2\pi}v^B(p+re^{\mathrm{i}t})\,\mathrm{d}t
	\]
	对充分小的坐标圆周成立。在盘内 \(v^B\) 调和，在盘外它等于 \(v\)，因此局部次平均检验覆盖全部点。最后，若闭包上的最大值大于边界最大值，它在内点取得，由强最大值原则与边界连续性矛盾。
\end{proof}

\begin{theorem}
	\label{g4:thm:0501:dirichlet}
	设 \(X\) 为黎曼曲面，\(\Omega\Subset X\) 为边界非空的光滑相对紧区域。对每个 \(f\in C(\partial\Omega,\mathbb R)\)，存在唯一
	\[
		u\in C(\overline\Omega,\mathbb R)\cap C^\infty(\Omega,\mathbb R),
		\qquad \Delta_g u=0,\quad u|_{\partial\Omega}=f.
	\]
	并且 \(\min_{\partial\Omega}f\le u\le\max_{\partial\Omega}f\)。
\end{theorem}
\begin{proof}
	设 \(m=\min f\)、\(M=\max f\)，令 \(\mathcal F\) 是所有在 \(\overline\Omega\) 上连续、在 \(\Omega\) 内次调和且边值不超过 \(f\) 的函数组成的族。常数 \(m\) 属于此族，而最大值原则使每个成员都不超过 \(M\)。定义
	\[
		u(x)=\sup_{v\in\mathcal F}v(x),\qquad x\in\Omega.
	\]
	该族对有限最大值及内盘 Poisson 替换封闭。构造的关键是证明上确界本身调和，并让它取到所需边值。

	固定坐标圆盘 \(B\Subset\Omega\) 及中心 \(a\)。选 \(v_n\in\mathcal F\) 使 \(v_n(a)\to u(a)\)，将它们依次替换为 \(\max(m,v_1,\ldots,v_n)\)，再作 \(B\) 内的 Poisson 替换，记所得调和部分为 \(h_n\)。它们单调增加且介于 \(m,M\) 之间。\cref{g4:lem:0501:poisson}的导数估计与\cref{g4:thm:0104:montel}证明中的紧空间选列论证，在内盘穷竭上逐次选子列，得到局部各阶一致收敛的子列；单调性使整个序列有同一个调和极限 \(h\)。这里按各阶导数再作对角选取，微分与极限相容，故极限满足调和方程。并有 \(h(a)=u(a)\)、\(h\le u\)。

	对任意 \(v\in\mathcal F\)，把 \(\max(v,m,v_1,\ldots,v_n)\) 作同样替换，得调和函数 \(k_n\ge h_n\)。在中心，\(0\le k_n(a)-h_n(a)\le u(a)-h_n(a)\to0\)。对非负调和函数 \(k_n-h_n\) 使用式\eqref{g4:eq:0501:harnack-disk}，知其在任一内盘上一致趋零。由于 \(v\le k_n\)，得到 \(v\le h\)。取所有 \(v\) 的上确界便有 \(u=h\)，所以 \(u\) 在 \(B\) 内调和。中心任意，故整个 \(\Omega\) 上调和。

	接着在每个边界点 \(p\) 构造障碍。取 \(z(p)=0\) 的全纯坐标，旋转实坐标，使边界为 \(y=\psi(x)\)，\(\Omega\) 在上侧，且 \(\psi(0)=\psi'(0)=0\)。光滑性给出 \(|\psi(x)|\le Cx^2\)。在下侧取圆心 \(a=-\mathrm{i}r\)、半径 \(r\) 的足够小圆，令 \(1/(2r)>C\)；其上半圆满足 \(y\le-x^2/(2r)\)，故除 \(p\) 外严格位于区域外。缩小坐标邻域后，
	\[
		b(z)=\log\frac{|z-a|}{r}
	\]
	在 \(\overline\Omega\) 的该邻域中非负、只在 \(p\) 为零，并在区域内调和。取围住 \(p\) 的较小坐标邻域，其人工边界与 \(\overline\Omega\) 的交集上有 \(b\ge c>0\)。在这个邻域内取 \(\min(b,c/2)\)，邻域外取常数 \(c/2\)，便得到全局连续超调和函数 \(\beta_p\)。两部分在人工边界附近已经相等；最小值保持超调和性。它在 \(p\) 为零，在 \(\overline\Omega\setminus\{p\}\) 上严格为正。

	给定 \(\varepsilon>0\)，由 \(f\) 连续，\(p\) 附近的边界上有 \(|f-f(p)|<\varepsilon\)。其余边界紧且 \(\beta_p\) 有正下界，因此可取 \(A>0\) 使
	\[
		f(p)-\varepsilon-A\beta_p\le f
		\le f(p)+\varepsilon+A\beta_p
		\quad\text{于 }\partial\Omega.
	\]
	左边是 \(\mathcal F\) 的成员。右边超调和；每个 \(v\in\mathcal F\) 与它作差为次调和函数，边界非正，所以内部也非正。取上确界得到相同的双侧夹估。令 \(x\to p\) 再令 \(\varepsilon\downarrow0\)，便有 \(u(x)\to f(p)\)。按这些值延拓，\(u\) 在闭包连续。唯一性及上下界由最大值原则得到。
\end{proof}

这就是\term[perronfangfa]{Perron 方法}[Perron method]\footnote{Oskar Perron（1880--1975），德国数学家，研究分析、连分数与代数。}：在所有满足边界不等式的次调和函数中取最大的候选，圆盘内的调和替换使候选达到调和性，边界障碍再保证它取到给定数据。光滑边界在此用于外切圆障碍；内部构造只需要局部坐标圆盘。

\subsection{Green 函数及对数奇点}
\label{g4:subsec:0501:03}

指定一个点 \(p\in\Omega\)，我们要将边界值为零的条件与 \(p\) 处的对数奇点同时满足。局部的 \(-\log|z|\) 不一定在整个曲面上有意义，因而不能直接拿它减去一个全局调和函数。下面先挖去一个小圆盘，在余下的光滑带边界区域上求解，再使小圆盘收缩。

\begin{theorem}
	\label{g4:thm:0501:green}
	设 \(\Omega\Subset X\) 为黎曼曲面中的光滑相对紧区域，\(\partial\Omega\ne\varnothing\)，且 \(p\in\Omega\)。存在唯一函数 \(G_\Omega(\cdot,p)\)，在 \(\overline\Omega\setminus\{p\}\) 上连续、在 \(\Omega\setminus\{p\}\) 上调和，边界值为零，并且在任一中心全纯坐标中有
	\begin{equation}
		\label{g4:eq:0501:green-local}
		G_\Omega(z,p)=-\log|z|+H(z),
	\end{equation}
	其中 \(H\) 在包含 \(z=0\) 的邻域内调和。它在 \(\Omega\setminus\{p\}\) 上严格为正，称为区域的\term[greenhanshu]{Green 函数}[Green function]。
\end{theorem}
\begin{proof}
	取中心坐标，使 \(\overline{D(0,2R)}\Subset\Omega\)。先为挖孔问题构造一个独立于孔半径的上界。在 \(\Omega\setminus\overline{D(0,R)}\) 上用\cref{g4:thm:0501:dirichlet}求解 \(h\)，其内圆周边值为 \(1\)，外边界值为 \(0\)。这个区域连通：原区域中的路径若穿过所删圆盘，可在圆盘外的一条薄圆环中把相应段绕行。强最大值原则给出 \(0<h<1\)。设 \(m_1=\max_{|z|=2R}h<1\)，比较圆环 \(R<|z|<2R\) 上的边值，得到
	\[
		h(z)\le1-a\log(|z|/R),\qquad
		a=\frac{1-m_1}{\log2}>0.
	\]
	此外 \(h\) 可光滑延拓过内圆周。确实，局部 Möbius 坐标把圆弧变为直线段，\(h-1\) 在该段连续取零；在一个小半盘中，把其半圆边值奇延拓到整圆周，Poisson 延拓也为奇函数。它在上半盘与原函数具有相同边值，由唯一性相等，故给出越过直径的调和延拓。搬回坐标即得所需光滑性。由上述圆环比较取右导数，有 \(\partial_rh|_{r=R}\le-a/R\)。

	取常数 \(A>1/a\)，并增大到 \(A+\log R\ge0\)，定义
	\[
		W(z)=
		\begin{cases}
			-\log(|z|/R)+A,&0<|z|\le R,\\
			Ah(z),&z\in\Omega\setminus D(0,R).
		\end{cases}
	\]
	它连续、非负、在外边界为零，并在接缝两侧调和。沿半径由内向外的导数跳跃满足
	\[
		\partial_rW|_{R+}-\partial_rW|_{R-}
		=A\partial_rh|_R+\frac1R<0.
	\]
	这保证 \(W\) 超调和。可以直接用比较来检验：从 \(W\) 减去一条调和比较函数，若出现负的内部最小值，落在某一光滑侧时，调和最小值原则会把它传到接缝或比较区域边界；落在接缝时，内侧径向导数应非正、外侧应非负，使导数跳跃非负，与上式矛盾。因此边界比较成立。

	对 \(0<\varepsilon<\min(R,1)\)，在 \(\Omega_\varepsilon=\Omega\setminus\overline{D(0,\varepsilon)}\) 上求 Dirichlet 解 \(G_\varepsilon\)，取内边值 \(-\log\varepsilon\)、外边值 \(0\)。最大值原则与 \(W\) 的比较给出
	\[
		0\le G_\varepsilon\le W.
	\]
	在任何避开 \(p\) 的紧子集附近，右端一致有界，Poisson 导数估计因而使任意趋零的 \(\varepsilon_n\) 都有一个子列，在各紧子集上连同各阶导数收敛到调和函数 \(G\)。这里用可数坐标盘穷竭与对角选取；\(X\) 的第二可数性保证可数化。

	在 \(\varepsilon<|z|<R\) 中，\(G_\varepsilon+\log|z|\) 调和，内边值为零，外边值介于 \(\log R\) 与 \(A+\log R\) 之间。最大值原则给出一个与 \(\varepsilon\) 无关的绝对值界。因此 \(G+\log|z|\) 在穿孔盘上有界且调和。说明其奇点可去：对一般有界穿孔盘调和函数 \(v\)，在某个半径 \(R_0\) 的圆周上作 Poisson 延拓 \(V\)。在 \(\delta<|z|<R_0\) 上，差 \(v-V\) 的外边值为零、内边值绝对值至多为常数 \(C\)，故
	\[
		|v(z)-V(z)|\le C\frac{\log(R_0/|z|)}{\log(R_0/\delta)}.
	\]
	固定 \(z\ne0\) 后令 \(\delta\downarrow0\)，得到 \(v=V\)。所以该有界函数在中心可去，延拓为调和函数，式\eqref{g4:eq:0501:green-local}成立。

	由于 \(0\le G\le W\)，而 \(W\) 连续至外边界且取零，\(G\) 也连续取零边值。奇点附近 \(G\to+\infty\)；若它在其他内点为零，强最小值原则使它在连通穿孔区域恒为零，与奇点矛盾，故严格为正。两个候选之差在 \(p\) 处可去，在闭包上连续，且边值为零，由最大值原则恒为零。这同时证明唯一性，说明选子列和辅助坐标没有改变所得函数。

	最后若 \(w=cz+O(z^2)\)、\(c\ne0\)，则 \(\log|w/z|\) 在中心附近为调和函数，因此改变中心坐标只把式\eqref{g4:eq:0501:green-local}中的正则部分改动一个调和项，奇点系数始终为 \(1\)。
\end{proof}

\begin{example}
	\label{g4:exm:0501:disk-green}
	在单位圆盘中，直接核对得到
	\[
		G_{\mathbb D}(z,a)=\log\left|\frac{1-\overline a z}{z-a}\right|,
		\qquad a\in\mathbb D.
	\]
	分子在圆盘中无零点，边界上分子与分母模相等；在 \(a\) 附近只有 \(-\log|z-a|\) 的奇点，故唯一性识别出此式。特别地 \(G_{\mathbb D}(z,0)=-\log|z|\)。若 \(\Omega_1\Subset\Omega_2\) 为两个上述区域且 \(p\in\Omega_1\)，则
	\[
		G_{\Omega_1}(x,p)\le G_{\Omega_2}(x,p),\qquad x\in\Omega_1\setminus\{p\}.
	\]
	两函数之差在 \(p\) 处可去，在 \(\partial\Omega_1\) 上右侧非负，应用最小值原则即得。扩大区域使接地边界远离极点，Green 函数随之增大。
\end{example}

\subsection{调和共轭、周期和能量}
\label{g4:subsec:0501:04}

\cref{g4:eq:0501:star-laplacian}把局部共轭的微分写成了全局闭形式。局部共轭不必能粘成单值函数；绕过洞或对数奇点时，其常数差可能累积为非零周期。这些周期将决定指数映射能否单值。

\begin{lemma}
	\label{g4:lem:0501:closed-form-periods}
	设 \(U\) 为光滑曲面区域，\(\alpha\) 为闭的光滑实一形式。\(\alpha\) 沿分片光滑道路的积分在定端点连续同伦下不变，沿闭路的积分给出 \(\pi_1(U,p)\to(\mathbb R,+)\) 的群同态。特别地，单连通区域上的闭一形式为恰当形式。
\end{lemma}
\begin{proof}
	先在坐标矩形中写 \(\alpha=a\,\mathrm{d}x+b\,\mathrm{d}y\)，闭性为 \(b_x=a_y\)。固定 \((x_0,y_0)\)，函数
	\[
		F(x,y)=\int_{x_0}^x a(s,y_0)\,\mathrm{d}s
		+\int_{y_0}^y b(x,t)\,\mathrm{d}t
	\]
	满足 \(F_x=a\)、\(F_y=b\)，故每点附近有原函数。两个原函数在交叠的每个连通分支上相差常数。

	对一条连续道路，按紧性作有限细分，使每段落在一个这样的原函数邻域内，以各段原函数的终值减初值之和定义积分。细分不改变总和；若同一段用了另一个原函数，其路径像是交叠中的连通集，差为常数，所以总和也不变。与通常的分片光滑积分比较，逐段基本积分公式说明两者相同。

	给定连续同伦 \(H:[0,1]^2\to U\)，原函数邻域的逆像覆盖参数方形。紧性使方形可分成有限小矩形，使每个小矩形的像位于一个原函数邻域内。每个小矩形边界的积分为零；对所有小矩形相加，内部公共边方向相反而相消。定端点同伦的两侧边为常值，故上下两条道路的积分相等。拼接使积分相加，反向使积分变号，于是得到所述群同态。单连通时所有闭路周期为零，沿基点到终点积分便定义全局原函数。
\end{proof}

\begin{proposition}
	\label{g4:prop:0501:conjugate-periods}
	设 \(u\) 是连通黎曼曲面区域 \(U\) 上的实调和函数。局部存在实函数 \(v\)，使 \(u+\mathrm{i}v\) 全纯。全局存在这样的单值 \(v\)，当且仅当
	\[
		\int_\gamma*\mathrm{d}u=0
	\]
	对 \(U\) 中所有分片光滑闭路成立；存在时 \(v\) 唯一到实常数。若所有周期只要求属于 \(2\pi\mathbb Z\)，则局部表达 \(\exp(u+\mathrm{i}v)\) 可选常数后粘成全局处处非零的全纯函数。
\end{proposition}
\begin{proof}
	局部构造已由\cref{g4:prop:0501:conformal-harmonicity}得到，且 Cauchy--Riemann 方程等价于 \(\mathrm{d}v=*\mathrm{d}u\)。若全局 \(v\) 存在，则闭路积分为零。反过来，固定基点，沿从基点到 \(x\) 的道路积分 \(*\mathrm{d}u\)。任意两条道路之差构成闭路，零周期使积分独立于道路；在局部坐标中改变终点，基本积分公式给出 \(\mathrm{d}v=*\mathrm{d}u\)。局部共轭之差为常数，故所得函数确为共轭，唯一性由连通性得出。

	若周期属于 \(2\pi\mathbb Z\)，道路积分的实值只确定到 \(2\pi\mathbb Z\)，但其指数 \(e^{\mathrm{i}v}\) 唯一。以局部原函数计算便知 \(e^{u+\mathrm{i}v}\) 全纯且不为零。这里整数周期允许共轭本身多值，同时让指数成为单值。
\end{proof}

对 \(u=\log|z|\)，有 \(*\mathrm{d}u=\mathrm{d}\arg z\)，逆时针小圆周上的周期为 \(2\pi\)。因而 \(e^{u+\mathrm{i}v}\) 可延拓为一个一阶零点。对 \(u=-G_\Omega(\cdot,p)\)，式\eqref{g4:eq:0501:green-local}也给出同样的局部周期 \(2\pi\)；区别只在正则的调和项。这正是下一节把 Green 函数转成圆盘映射时需要的奇点行为。

\begin{proposition}
	\label{g4:prop:0501:energy}
	设 \(\Omega\Subset X\) 为光滑相对紧区域，\(g\) 为光滑共形度量，\(u,v\) 为 \(\overline\Omega\) 上的实光滑函数。则
	\begin{equation}
		\label{g4:eq:0501:green-identity}
		\int_\Omega\langle\mathrm{d}u,\mathrm{d}v\rangle_g\,\mathrm{d}A_g
		=\int_\Omega u\Delta_gv\,\mathrm{d}A_g
		+\int_{\partial\Omega}u\,\partial_\nu v\,\mathrm{d}s_g,
	\end{equation}
	其中 \(\nu\) 为外单位法向量。能量
	\[
		E_\Omega(u)=\int_\Omega|\mathrm{d}u|_g^2\,\mathrm{d}A_g
	\]
	不受共形度量的选择影响。若 \(u\) 调和，则它在所有与自己有相同边值的光滑函数中唯一使能量最小。
\end{proposition}
\begin{proof}
	坐标计算给出
	\[
		\mathrm{d}u\wedge*\mathrm{d}v
		=(u_xv_x+u_yv_y)\,\mathrm{d}x\wedge\mathrm{d}y
		=\langle\mathrm{d}u,\mathrm{d}v\rangle_g\,\mathrm{d}A_g.
	\]
	共形因子的倒平方与面积的平方因子相消，得到能量不变性。再由乘积法则及式\eqref{g4:eq:0501:star-laplacian}，
	\[
		\mathrm{d}(u*\mathrm{d}v)
		=\mathrm{d}u\wedge*\mathrm{d}v-u\Delta_gv\,\mathrm{d}A_g.
	\]
	第一卷第14.4节定理14.4.3的 Stokes 公式采用外法向优先的边界定向。在正向边界单位切向量 \(T\) 上，\((*\mathrm{d}v)(T)=\partial_\nu v\)，所以积分得到式\eqref{g4:eq:0501:green-identity}。

	若 \(w=u+\phi\) 与 \(u\) 边值相同，则 \(\phi|_{\partial\Omega}=0\)。将式\eqref{g4:eq:0501:green-identity}用于 \((\phi,u)\)，由 \(\Delta_gu=0\) 知交叉项为零，于是
	\[
		E_\Omega(w)=E_\Omega(u)+E_\Omega(\phi)\ge E_\Omega(u).
	\]
	等号迫使 \(\mathrm{d}\phi=0\)，连通性使 \(\phi\) 为常数，零边值又使该常数为零。这里极小性质是在已经构造出调和解后证明的，边值只有连续性时并未预设能量有限。
\end{proof}

\begin{example}
	\label{g4:exm:0501:annulus-energy}
	在 \(A(r,R)=\{r<|z|<R\}\) 上，把内边界值取为 \(0\)、外边界值取为 \(1\)，解为
	\[
		u(z)=\frac{\log(|z|/r)}{\log(R/r)}.
	\]
	其共轭微分沿正向同心圆的周期是 \(2\pi/\log(R/r)\)，能量为
	\[
		E_{A(r,R)}(u)
		=\int_r^R\int_0^{2\pi}
		\frac{1}{s^2\log(R/r)^2}s\,\mathrm{d}\theta\,\mathrm{d}s
		=\frac{2\pi}{\log(R/r)}.
	\]
	因此环越长，连接两条边界的单位电势差所需能量越小；同一个对数比值还将在极值长度中出现。
\end{example}

\subsection{最大值原理与 Harnack 紧性}
\label{g4:subsec:0501:05}

圆盘内的平均值与核估计可以沿曲面传播，形成 Harnack 不等式。\footnote{Karl Gustav Axel Harnack（1851--1888），德国数学家，研究势论与实代数曲线。}关键是用有限条相交坐标盘连接一个基点和所考察的紧集：每经过一张图，就把一个点的控制传给下一个点。这个过程只耗费有限个常数，所以保留了适于极限的统一估计。

\begin{proposition}
	\label{g4:prop:0501:harnack-compactness}
	设 \(U\) 为连通黎曼曲面区域，\(p\in U\)。实调和函数若在 \(U\) 内取得局部最大值或局部最小值，则为常数。对每个紧集 \(K\Subset U\)，存在 \(C_K\ge1\)，使所有非负调和函数满足
	\[
		C_K^{-1}h(p)\le h(x)\le C_Kh(p),\qquad x\in K.
	\]
	若非负调和函数列 \(h_n\) 的 \(h_n(p)\) 有界，则有子列在所有紧子集上连同各阶导数收敛到非负调和函数。若在 \(p\) 的值趋于正无穷，则 \(h_n\) 在每个紧子集上一致趋于正无穷。实调和函数列若在各紧子集上绝对值一致有界，也具有上述各阶局部收敛子列。
\end{proposition}
\begin{proof}
	局部极值先在一张坐标图中由\cref{g4:cor:0103:harmonic-maximum}使函数局部为常数。任意道路都可由有限个相交坐标盘覆盖；在相邻交叠中已有一个开集取常值，全纯导数的恒等定理使常值传遍下一盘。因此它在整个连通区域为常数。

	对每个 \(x\in K\)，选一条连接 \(p,x\) 的道路，用有限个相交坐标圆盘覆盖，并把道路上的连接点放在各圆盘的较小同心盘内。式\eqref{g4:eq:0501:harnack-disk}在相邻连接点间给出双侧常数比较，有限次相乘，得到 \(h(p)\) 与 \(x\) 附近各点值的比较。由紧性取有限个这样的终点邻域覆盖 \(K\)，再取各常数的最大值，便得到 \(C_K\)。这称为\term[harnackbudengshi]{Harnack 不等式}[Harnack inequality]在曲面上的形式。

	若 \(h_n(p)\) 有界，前一结论给出每个紧集上的函数一致界。对包含 \(K\) 的有限组稍大坐标盘使用 Poisson 导数估计，得到各阶导数的一致界及等度连续性。对紧集穷竭、有限图册和导数阶数作对角选取，\cref{g4:thm:0104:montel}证明中的紧空间选列论证给出一个各阶局部一致收敛的子列。二阶导数收敛使极限满足调和方程，函数值收敛保留非负性。若 \(h_n(p)\to\infty\)，则 \(h_n(x)\ge C_K^{-1}h_n(p)\) 给出紧集上一致发散。最后，实调和函数列已有局部绝对值界时，可直接从导数估计开始作同一紧性论证。
\end{proof}

特别地，一个单调增加的非负调和函数列要么在所有紧子集上收敛到调和函数，要么处处并局部一致地趋于正无穷：只须观察它在一个基点的值有界还是无界。有限分支中的唯一逐点极限使全部收敛子列相同，导数估计继而保证整列局部收敛。统一化证明将对不断增大的曲面区域使用这类控制；正性、一个基点的归一化及紧集上的局部估计，分别承担传播上界、阻止退化与提取极限的作用。
% 纲要条目排为编号小节标题，并进入目录与书签；说明段仍保留于纲要注释中。
% !TeX root = ../../Book4.tex
% 纲要标题：### 5.2 带边界单连通区域的映照
% 纲要位置：几何学计划纲要.md，第 3641 行。
% 正文状态：初稿；尚未正式审核。

% 纲要内容开始
% 
% * 拓扑圆盘和边界正则化
% * 由 Green 函数构造全纯映射
% * 边界模长与映射次数
% * 单叶性和圆盘映照
% 
% 纲要内容结束

\section{带边界单连通区域的映照}
\label{g4:sec:0502}

在一个没有预先给定平面坐标的曲面区域上，Green 函数仍能测量点到边界的势差。它在极点附近像 \(-\log|z|\)，而边界值为零。如果把它补成局部全纯函数，再取负指数，极点就变成一个简单零点，边界则变成模长为 \(1\) 的限制。下面证明这个构造给出整个单位圆盘的单叶坐标。单连通性用于消除指数的多值性；边界条件用于保证原像不逃向边界。

\subsection{拓扑圆盘和边界正则化}
\label{g4:subsec:0502:01}

设 \(X\) 为黎曼曲面。称紧子集 \(D\subset X\) 为嵌入的闭拓扑圆盘，是指存在从闭单位圆盘到 \(D\) 的同胚，且它把开圆盘送到 \(D\) 在 \(X\) 中的内部。以后写 \(U=\operatorname{int}D\)。若 \(\partial D\) 是光滑嵌入圆周，则 \(D\) 同时是紧光滑带边界曲面，\(U\Subset X\)。这里的拓扑条件只描述区域的形状，并没有提供全纯坐标；把它变成解析圆盘正是本节的目标。

\ProseParagraphBreak

由有限条坐标弧围成的盘常带有角点。势论中的光滑边界条件可以通过局部修改实现，并且修改能够避开事先选定的内部紧集。下述引理只处理已经得到的拓扑盘；下一节另行证明一般非紧单连通曲面如何被这样的盘穷竭。

\begin{lemma}
	\label{g4:lem:0502:round-corners}
	设 \(D\) 是黎曼曲面 \(X\) 中的嵌入闭拓扑圆盘，其边界为分片光滑正则曲线，只有有限个角点，且每个角点处的两条单侧切线横截。给定紧集 \(K\subset\operatorname{int}D\) 及包含 \(D\) 的开集 \(W\)，存在光滑边界的嵌入闭拓扑圆盘 \(D'\subset W\)，满足 \(K\subset\operatorname{int}D'\)。修改可以限制在各角点任意预先选定的充分小邻域内。
\end{lemma}
\begin{proof}
	取彼此不交、避开 \(K\) 且闭包包含于 \(W\) 的角点坐标邻域。只需说明一个角点附近的构造，因为随后可以在这些不交邻域中分别实施。

	将角点坐标平面中的坐标轴作实线性变换后，两条单侧弧都能以横坐标为参数，并分别位于横坐标的两侧。因此在一个小矩形中，边界写成 \(y=b(x)\)：函数 \(b\) 连续、分段光滑，在 \(x=0\) 处只有一阶导数的跳跃；区域在这张图像的一侧。这样的坐标方向存在，因为两条离开角点的切向射线不同且不共线，可以选择一个线性泛函，使它在两条射线上取相反符号。缩小矩形后，两弧在这个方向上的导数都不为零，便得到所需图像。

	把 \(b\) 在稍大的区间上延拓为连续函数，取光滑非负、积分为 \(1\) 的紧支核，作卷积 \(b*\varphi_\varepsilon\)。卷积在紧区间上一致逼近 \(b\)。选光滑截断函数 \(\chi(x)\)，在 \(0\) 附近为 \(1\)，在上述小区间端部为 \(0\)，令
	\[
		b_\varepsilon(x)=\chi(x)(b*\varphi_\varepsilon)(x)
		+(1-\chi(x))b(x).
	\]
	这是光滑函数：在唯一可能不光滑的点 \(0\) 附近只剩卷积项；在截断过渡区，原函数本来光滑。它与 \(b\) 在端部相等，并一致趋于 \(b\)。

	还需保证换图像不会改变盘的拓扑。记 \(\delta(x)=b_\varepsilon(x)-b(x)\)，取一个光滑函数 \(\beta\)，在 \(0\) 附近为 \(1\)，支集放在很短的纵向区间内。在矩形中定义
	\[
		T(x,y)=\bigl(x,y+\beta(y-b(x))\delta(x)\bigr).
	\]
	取 \(\varepsilon\) 充分小，使 \(\|\beta'\|_\infty\|\delta\|_\infty<1/2\)。对每个固定 \(x\)，第二分量关于 \(y\) 严格递增，并在区间端部等于 \(y\)，因而是该纵向区间到自身的同胚。它的斜率至少为 \(1/2\)，所以逆映射连续依赖 \(x,y\)。于是 \(T\) 是矩形的同胚，在边界附近恒等，且把旧图像送到新图像。将它在矩形外延为恒等映射，便得到 \(X\) 的一个同胚，其像盘的边界在该角点处已经光滑。

	对所有角点作上述修改。所得环境同胚在 \(K\) 附近恒等，修改支集包含在 \(W\) 内，故 \(D'=T(D)\) 仍为闭拓扑盘，包含 \(K\) 于内部且包含于 \(W\)。边界在修改区为光滑图像，在接缝附近与旧光滑弧完全相同，所以全边界光滑。
\end{proof}

这个构造改变的是区域边界，曲面已有的复结构始终保留。因而可以在正则化后的区域直接使用\cref{g4:thm:0501:dirichlet,g4:thm:0501:green}。只要极点与需要控制的紧集留在内部，后面的映照构造就适用。

\subsection{由 Green 函数构造全纯映射}
\label{g4:subsec:0502:02}

现在固定光滑边界闭拓扑盘 \(D\subset X\)，令 \(U=\operatorname{int}D\)，并选 \(p\in U\)。取 \(p\) 附近的全纯坐标 \(z\)，满足 \(z(p)=0\)。上一节构造了 Green 函数 \(G=G_U(\cdot,p)\)，它在 \(U\setminus\{p\}\) 调和且为正，连续延至 \(\partial U\) 取零值，并在 \(p\) 附近满足
\begin{equation}
	\label{g4:eq:0502:green-local}
	G=-\log|z|+h,
\end{equation}
其中 \(h\) 在整个小坐标盘内调和。

\ProseParagraphBreak

依照上一节的约定，\(*\mathrm{d}x=\mathrm{d}y\)、\(*\mathrm{d}y=-\mathrm{d}x\)。因为 \(G\) 调和，\(*\mathrm{d}G\) 在去点区域上是闭一形式。它的局部原函数 \(v\) 满足 \(\mathrm{d}v=*\mathrm{d}G\)，所以 \(G+\mathrm{i}v\) 局部全纯。绕极点一周时，\(v\) 不会回到原值；需要先算出这个变化，再判断指数能否成为单值函数。

\begin{lemma}
	\label{g4:lem:0502:green-periods}
	设 \(U\) 是光滑边界闭拓扑盘的内部，\(p\in U\)，\(G=G_U(\cdot,p)\)。对每条分片光滑闭路 \(\gamma\subset U\setminus\{p\}\)，有
	\[
		\int_\gamma *\mathrm{d}G\in2\pi\mathbb Z.
	\]
	若 \(\gamma\) 是以 \(p\) 为中心的全纯坐标中的小圆周，取逆时针方向，则该积分为 \(-2\pi\)。
\end{lemma}
\begin{proof}
	先计算小圆周。由式\eqref{g4:eq:0502:green-local}，
	\[
		*\mathrm{d}G=-\mathrm{d}\arg z+*\mathrm{d}h.
	\]
	\(h\) 在整个坐标盘内调和，故在该盘内有单值调和共轭，后一项是恰当形式。沿逆时针小圆周积分，结果为 \(-2\pi\)。

	说明这个局部计算为何控制全部闭路。以 \(p\) 的一个小坐标圆盘 \(B\) 覆盖缺失点，便有开覆盖
	\[
		U=(U\setminus\{p\})\cup B,
		(U\setminus\{p\})\cap B=B\setminus\{p\}.
	\]
	这三个开集均道路连通。\(B\setminus\{p\}\) 径向形变收缩到小圆周，其基本群由该圆周生成，这是第一卷第5.1节定理5.1.13的圆周绕数计算。第一卷第5.2节定理5.2.4的 van Kampen 定理\footnote{Egbert Rudolf van Kampen（1908--1942），生于比利时的数学家，在荷兰求学、美国任教，研究拓扑与拓扑群。}说明：因为 \(B\) 和 \(U\) 都单连通，\(\pi_1(U\setminus\{p\})\) 由这条小圆周的共轭及其逆元正规生成。

	由\cref{g4:lem:0501:closed-form-periods}，闭一形式的积分只依赖定端点连续同伦类，并对道路拼接可加；因此它给出从基本群到实数加法群的同态。沿共轭道路 \(\alpha*\gamma*\bar\alpha\) 的积分等于沿 \(\gamma\) 的积分，连接弧的贡献相消。所有闭路的周期于是都是小圆周周期 \(-2\pi\) 的整数倍。
\end{proof}

\begin{proposition}
	\label{g4:prop:0502:green-map}
	设 \(D\) 是黎曼曲面中的光滑边界闭拓扑盘，\(U=\operatorname{int}D\)，\(p\in U\)。存在全纯函数 \(F:U\to\mathbb D\)，满足
	\[
		F(p)=0,\qquad |F(q)|=e^{-G_U(q,p)}\quad(q\ne p).
	\]
	它只在 \(p\) 有零点，且该零点为简单零点；这样的函数彼此只差一个模为 \(1\) 的常数因子。若在中心坐标 \(z\) 中写 \(G=-\log|z|+h\)，则
	\[
		\left|\frac{\mathrm{d}F}{\mathrm{d}z}(p)\right|=e^{-h(p)}.
	\]
\end{proposition}
\begin{proof}
	在 \(U\setminus\{p\}\) 中选基点 \(q_0\)，先给局部共轭指定一个实常数。对从 \(q_0\) 到 \(q\) 的道路积分 \(*\mathrm{d}G\)，得到一个可能依赖道路的数 \(v(q)\)。由\cref{g4:lem:0502:green-periods}，改变道路使这个数增加 \(2\pi\) 的整数倍。因此
	\begin{equation}
		\label{g4:eq:0502:exponential-map}
		F(q)=\exp\bigl(-G(q)-\mathrm{i}v(q)\bigr)
	\end{equation}
	与道路无关。在任一不含 \(p\) 的小坐标盘上，\(G+\mathrm{i}v\) 是全纯函数，因而 \(F\) 全纯且不为零。\(G>0\) 给出 \(|F|<1\)。

	在 \(p\) 附近，给 \(h\) 选单值调和共轭 \(h^*\)，使 \(H=h+\mathrm{i}h^*\) 全纯。在一个割开的小盘中，\(v=-\arg z+h^*+c\)，其中 \(c\) 为实常数。式\eqref{g4:eq:0502:exponential-map}便成为
	\[
		F=e^{-\mathrm{i}c}\,z\,e^{-H(z)}.
	\]
	右边在完整小盘上全纯，因此补上 \(F(p)=0\) 后得到 \(U\) 上的全纯函数；该式还说明零点恰为一阶，且导数模为 \(e^{-h(p)}\)。

	若 \(F_1,F_2\) 都有上述模长，二者在 \(p\) 的简单零点相消，故 \(F_1/F_2\) 在 \(U\) 上处处全纯，模恒为 \(1\)。最大模原理使这个比值为常数。反之，乘任一模为 \(1\) 的常数显然保留全部要求。
\end{proof}

\subsection{边界模长与映射次数}
\label{g4:subsec:0502:03}

Green 边界值给出 \(|F|\to1\)，这已经足以控制所有位于圆盘内部的原像。确实，\(|F|\) 在 \(\overline U\) 上连续延拓，边界值为 \(1\)。对任意 \(0<r<1\)，集合
\[
	\{q\in U:|F(q)|\le r\}
\]
是 \(\overline U\) 中避开边界的闭集，因此紧且包含在 \(U\) 内。目标中的每个紧集都包含在某个 \(\{|w|\le r\}\) 中，其原像便紧。也就是说，\(F:U\to\mathbb D\) 是固有映射。这里控制的是模长；证明不要求先把 \(F\) 的相位逐点延拓到边界。

\ProseParagraphBreak

第二章的次数定理以紧源曲面为假设。当前的 \(U\) 是开曲面，所需的紧性由固有性提供。下面将对应的计数论证写出，使这个变化的作用明确可见。

\begin{proposition}
	\label{g4:prop:0502:proper-degree}
	设 \(S,T\) 是连通黎曼曲面，\(f:S\to T\) 是固有的非恒定全纯映射。则 \(f\) 满射，每个纤维有限，且存在整数 \(d\ge1\)，使对每个 \(y\in T\) 都有
	\[
		\sum_{x\in f^{-1}(y)}e_x(f)=d.
	\]
	这个整数称为 \(f\) 的次数。若某个纤维只有一个点，且该点的分歧指数为 \(1\)，则 \(f\) 为双全纯映射。
\end{proposition}
\begin{proof}
	由\cref{g4:thm:0202:local-model}，\(f\) 的局部形式为 \(z\mapsto z^e\)，所以它开，纤维中的各点也孤立。固有性使纤维紧；紧离散空间有限，所以每个纤维有限。

	像还在 \(T\) 中闭。若像中的点列 \(y_n\) 趋于 \(y\)，可以把尾部与 \(y\) 放在同一个紧坐标闭盘内。任选 \(x_n\in f^{-1}(y_n)\)，固有性使这些原像位于同一个紧集。取收敛子列 \(x_{n_k}\to x\)，连续性给出 \(f(x)=y\)。黎曼曲面有可数局部基，这个点列判据足以证明闭性。非空的像既开又闭，而 \(T\) 连通，故 \(f\) 满射。

	固定 \(y\)，记纤维为 \(x_1,\ldots,x_m\)。在各点取闭包互不相交的小坐标盘 \(B_i\)，使 \(f\) 的局部幂模型成立，并使 \(x_i\) 是 \(\overline B_i\) 中 \(y\) 的唯一原像。充分接近 \(y\) 的目标点，其全部原像都在这些盘内：否则可取 \(y_n\to y\) 和 \(x_n\notin\bigcup_iB_i\) 满足 \(f(x_n)=y_n\)，固有性再次给出收敛子列，其极限既在这个开集的补集中，又属于 \(f^{-1}(y)\)，矛盾。

	目标坐标中取充分小的 \(y'\) 邻域后，局部幂模型说明 \(B_i\) 内全部原像的重数之和恒为 \(e_{x_i}(f)\)；对 \(y'\ne y\)，这些原像实际上互异且各重数为 \(1\)。因此总重数关于 \(y\) 局部常值。\(T\) 连通，得到固定的正整数 \(d\)。若某个纤维总重数为 \(1\)，则每个纤维都恰有一个点，而且处处分歧指数为 \(1\)。于是 \(f\) 双射，局部全纯逆又按唯一性拼成全局全纯逆。
\end{proof}

对 Green 函数构造的 \(F\)，零纤维恰为 \(\{p\}\)，且刚已证明 \(e_p(F)=1\)。所以它的次数为 \(1\)。这一步同时得到满射性和单射性：边界条件保证每个值都能取到，而一个简单零点把所有纤维的总重数固定为 \(1\)。

\subsection{单叶性和圆盘映照}
\label{g4:subsec:0502:04}

把上述构造与计数合在一起，就得到适合在穷竭中反复使用的归一化映照。归一化必须说明在哪张源坐标中测量导数，因为曲面上两个不同的坐标会给出不同的导数数值。

\begin{theorem}
	\label{g4:thm:0502:bordered-disk}
	设 \(X\) 为黎曼曲面，\(D\subset X\) 是光滑边界的嵌入闭拓扑圆盘，\(U=\operatorname{int}D\)。给定 \(p\in U\) 及 \(p\) 附近的中心全纯坐标 \(z\)，存在唯一双全纯映射
	\[
		f_U:U\longrightarrow\mathbb D,
		\qquad f_U(p)=0,\qquad
		\frac{\mathrm{d}f_U}{\mathrm{d}z}(p)>0.
	\]
	最后一个不等式指正实数。它满足
	\[
		G_U(q,p)=-\log|f_U(q)|\quad(q\ne p).
	\]
	若 \(G_U=-\log|z|+h\) 于 \(p\) 附近，则 \(f_U'(p)=e^{-h(p)}\)。
\end{theorem}
\begin{proof}
	由\cref{g4:prop:0502:green-map}构造 \(F\)。边界模长使它固有，\cref{g4:prop:0502:proper-degree}与简单零点使它双全纯。乘以唯一的模为 \(1\) 的常数，使 \(F\) 在坐标 \(z\) 中的非零导数成为正实数，便得所求 \(f_U\)，Green 公式与导数公式也随构造成立。

	若 \(f_1,f_2\) 都满足归一化条件，则 \(f_2\circ f_1^{-1}\) 是固定 \(0\) 的圆盘自同构。由\cref{g4:thm:0103:disk-aut}，它为旋转 \(w\mapsto e^{\mathrm{i}\theta}w\)。在基点求导，有 \(e^{\mathrm{i}\theta}=f_2'(p)/f_1'(p)>0\)，从而 \(e^{\mathrm{i}\theta}=1\)，两映射相同。
\end{proof}

\begin{corollary}
	\label{g4:cor:0502:domain-monotonicity}
	设 \(U\subset V\) 是同一黎曼曲面中两个光滑边界闭拓扑盘的内部，\(p\in U\)。在同一个中心坐标中取上述归一化映照 \(f_U,f_V\)。则
	\[
		0<f_V'(p)\le f_U'(p),\qquad
		G_V(q,p)\ge G_U(q,p)\quad(q\in U\setminus\{p\}).
	\]
	导数等号成立当且仅当 \(U=V\)。
\end{corollary}
\begin{proof}
	复合 \(H=f_V\circ f_U^{-1}:\mathbb D\to\mathbb D\) 全纯且固定 \(0\)。Schwarz 引理给出 \(|H(w)|\le|w|\) 及 \(|H'(0)|\le1\)。代入 \(w=f_U(q)\)，再用链式法则，得到
	\[
		|f_V(q)|\le|f_U(q)|,\qquad
		H'(0)=\frac{f_V'(p)}{f_U'(p)}\in(0,1].
	\]
	取负对数便得到 Green 函数的不等式。若导数比为 \(1\)，Schwarz 引理的等号条件使 \(H\) 为恒等映射，故 \(f_V(U)=\mathbb D=f_V(V)\)。\(f_V\) 单射，于是 \(U=V\)。反向由归一化映照的唯一性成立。
\end{proof}

\begin{example}
	\label{g4:exm:0502:disk-green}
	对平面圆盘 \(U=\{|z|<R\}\)，极点 \(p\in U\)，以 \(z-p\) 为中心坐标，归一化映照和 Green 函数分别为
	\[
		f_U(z)=\frac{R(z-p)}{R^2-\overline pz},\qquad
		G_U(z,p)=\log\left|\frac{R^2-\overline pz}{R(z-p)}\right|.
	\]
	分母在 \(U\) 内不为零，且
	\[
		f_U'(p)=\frac{R}{R^2-|p|^2}>0.
	\]
	当 \(p=0\) 时，\(G_U=\log(R/|z|)\)，构造所得坐标正是 \(z/R\)。这些公式依赖区域的共形结构；为定义 Laplace 算子而选的辅助共形度量不出现在答案中。
\end{example}

现在每个光滑相对紧拓扑盘都有唯一的归一化全纯坐标，嵌套区域上的基点导数也有确定的单调方向。要将这个局部于区域的结论推广到整张非紧曲面，还需构造合适的盘穷竭，并证明归一化映照的极限不会退化；下一节分别处理这两个问题。
% 纲要条目排为编号小节标题，并进入目录与书签；说明段仍保留于纲要注释中。
% !TeX root = ../../Book4.tex
% 纲要标题：### 5.3 一般单连通曲面的统一化
% 纲要位置：几何学计划纲要.md，第 3648 行。
% 正文状态：初稿；尚未正式审核。

% 纲要内容开始
% 
% * 非紧单连通曲面的圆盘穷竭
% 
%   利用相对紧曲面区域、边界曲线的填洞及单连通性建立圆盘穷竭；本引理在本章证明，不假定已掌握一般非紧曲面分类。
% 
% * 穷竭区域的归一化映照
% * 正规族极限及非退化性
% 
%   单列归一化导数、单叶函数畸变控制以及逆映射紧性的引理；分别排除常值极限并验证极限像，不以“由正规族”略过核心存在步骤。
% 
% * 双曲型与抛物型的区分
% * 紧单连通情形和黎曼球面
% * 万有覆盖与曲面的商表示
% 
% 纲要内容结束

\section{一般单连通曲面的统一化}
\label{g4:sec:0503}

一张单连通非紧曲面可能需要无穷多张坐标图。上一节只能处理其中一个相对紧圆盘；要得到整张曲面的坐标，必须让这些圆盘逐渐扩大，同时控制映照的尺度。我们先证明所需圆盘确实存在，再把尺度固定在同一点的一阶导数上。极限目标是有限半径圆盘还是整个平面，将由这些圆盘的共形半径决定。

\subsection{非紧单连通曲面的圆盘穷竭}
\label{g4:subsec:0503:01}

这里所说的光滑闭盘，是曲面中具有光滑边界、同胚于闭单位圆盘的紧二维子流形；这个定义涉及拓扑类型和已有的光滑结构，没有预先选择全局复坐标。下述论证只使用紧曲面的分类。

\begin{lemma}
	\label{g4:lem:0503:compact-neighborhood}
	设 \(X\) 是连通非紧黎曼曲面，\(K\subset X\) 紧。存在具有非空光滑边界的紧连通曲面区域 \(M\subset X\)，使 \(K\subset\operatorname{int}M\)。可以递归选取这类区域，使其内部穷尽 \(X\)。
\end{lemma}
\begin{proof}
	先用有限个相对紧坐标圆盘覆盖 \(K\)，并以有限条道路连接这些圆盘。所得紧连通集合记为 \(L\)，它包含 \(K\) 的一个邻域。第一卷第12.3节推论12.3.6的紧支截断给出光滑 \(\chi:X\to[0,1]\)，在 \(L\) 附近等于 \(1\)，且支集紧。第二卷第14.2节定理14.2.4的 Sard 定理\footnote{Arthur Sard（1909--1980），美国数学家，研究临界值、逼近与样条。}保证存在正则值 \(t\in(0,1)\)。取 \(\{\chi\ge t\}\) 中包含 \(L\) 的连通分支为 \(M\)。隐函数定理使其边界光滑；正则水平的紧性使边界为有限条圆周。它在紧支集内闭，故紧。若边界为空，\(M\) 在连通 \(X\) 中既开又闭，便有 \(M=X\)，与非紧性矛盾。

	第二可数性给出可数个相对紧坐标圆盘覆盖。第 \(n\) 步把前 \(n\) 个闭子圆盘和上一步区域一起放入 \(K\)，重复构造，便得到内部递增且覆盖 \(X\) 的区域。这一构造不要求 \(X\) 单连通。
\end{proof}

填洞之前，需要说明在曲面上能够填入的究竟是什么。平面上的 Jordan 定理不能直接用于尚无平面坐标的整张曲面。

\begin{lemma}
	\label{g4:lem:0503:null-circle}
	设 \(X\) 是单连通非紧黎曼曲面。每条光滑嵌入圆周 \(C\subset X\) 都恰好围出一个紧闭盘，其边界为 \(C\)。
\end{lemma}
\begin{proof}
	取 \(C\) 的零伦及其紧像，再用\cref{g4:lem:0503:compact-neighborhood}把这个像放在紧连通光滑曲面 \(M\) 的内部。因此 \(C\) 已在 \(M\) 中零伦，而 \(\partial M\ne\varnothing\)。\(M\) 由复结构定向。我们说明紧曲面中这一条件如何迫使 \(C\) 围盘。

	第一卷第6.6节的切开构造表明：若 \(C\) 不分离 \(M\)，就有一条闭曲线与它横截相交一次。这个模 \(2\) 交数在零伦下不变，零伦圆周的交数应为零，矛盾。此处用的是第二卷第16.2节定理16.2.8及推论16.2.9的同调交数；也可将紧零伦所在区域封盘后应用其闭曲面版本。因此 \(C\) 分开 \(M\)。

	选 \(M\) 的一条边界 \(b_0\)，记不含 \(b_0\) 的切开分支为 \(M_1\)，另一分支为 \(M_2\)。第一卷第6.4节定理6.4.10的带边界标准模型，使每一边都能用手柄环路和边界环路切成多边形。用连接弧把基点接到各边界，再按第一卷第5.2节的 van Kampen 计算，其表示为
	\[
		\left\langle a_1,b_1,\ldots,a_h,b_h,c_0,c_1,\ldots,c_{b-1}\ \middle|\
		\prod_{j=1}^{h}[a_j,b_j]c_0c_1\cdots c_{b-1}=1\right\rangle.
	\]
	这里 \(c_0\) 对应预留的边界 \(b_0\)，\(b=|\pi_0(\partial M)|\)。消去只出现一次的 \(c_0\)，余下字母就是自由基。为使这组基适应 \(C\)，先在 \(M_1,M_2\) 各自切开，再沿 \(C\) 粘合；各手柄字母与原边界字母分别来自两侧。在这个自由基中，\(C\) 的类共轭于 \(M_1\) 的边界字的逆，即
	\[
		\prod_{j=1}^{h_1}[\alpha_j,\beta_j]d_1\cdots d_k
	\]
	或其逆，其中 \(d_1,\ldots,d_k\) 是落在 \(M_1\) 的原边界，且都不是被消去的 \(c_0\)。不同字母独立，这个字自由约化后非空，除非 \(h_1=k=0\)。由于 \(C\) 零伦，必须出现后一情形；标准模型于是给出 \(M_1\cong\Sigma_{0,1}\)，即闭盘。该盘作为 \(X\) 的子集已有光滑边界 \(C\)。

	这个盘的内部是 \(X\setminus C\) 的一个分支。其外侧连通：若有道路经过盘内，可沿盘的边界领圈绕开；局部领圈也说明没有另外的补分支。若外侧亦为紧闭盘，两盘沿边界粘成 \(X\)，则 \(X\) 为紧球面，矛盾。所以只有一个紧闭盘侧。
\end{proof}

\begin{theorem}
	\label{g4:thm:0503:disk-exhaustion}
	设 \(X\) 是单连通非紧黎曼曲面，\(p\in X\)。存在光滑闭盘 \(\overline U_n\subset X\)，满足
	\[
		p\in U_1,\qquad \overline U_n\subset U_{n+1},\qquad
		X=\bigcup_{n\ge1}U_n.
	\]
	特别地，\(X\) 的每个紧集都包含在某个 \(U_n\) 中。
\end{theorem}
\begin{proof}
	先证任意紧集 \(K\) 均包含在一个光滑闭盘的内部。取上一引理中的紧连通区域 \(M\)，使 \(K\subset\operatorname{int}M\)。其边界只有有限条圆周 \(C_1,\ldots,C_b\)，每条都围出唯一紧盘 \(D_i\)。\(\operatorname{int}M\) 连通且不碰 \(C_i\)，故或者位于 \(D_i\) 内，或者位于其外侧。

	如果某个 \(D_i\) 包含 \(\operatorname{int}M\)，那么 \(K\subset\operatorname{int}D_i\)，任务完成。否则所有 \(D_i\) 都在 \(M\) 的外侧。这时它们的内部两两不交：两条边界互不相交，其盘若相交，只能有一个包含另一条边界；但每条 \(C_j\) 的内侧都接触 \(\operatorname{int}M\)，不能落在与 \(\operatorname{int}M\) 不交的另一盘的内部。于是将全部 \(D_i\) 补到 \(M\) 上，就得到一个无边界紧曲面子集。边界领圈保证这个子集在 \(X\) 中开，紧性保证它闭；连通性遂使它等于非紧的 \(X\)，矛盾。因此必有一个盘包含 \(K\)。

	取\cref{g4:lem:0503:compact-neighborhood}给出的紧区域穷竭 \(K_n\)，并将 \(p\) 加入 \(K_1\)。先用上述结论包住 \(K_1\)，随后每次包住 \(K_{n+1}\cup\overline U_n\)。所得闭盘满足全部要求。紧集被某个 \(U_n\) 包含，来自递增开覆盖的有限子覆盖性质。
\end{proof}

\subsection{穷竭区域的归一化映照}
\label{g4:subsec:0503:02}

固定 \(p\) 附近的局部坐标 \(z\)，令 \(z(p)=0\)。由\cref{g4:thm:0502:bordered-disk}，每个 \(U_n\) 有唯一双全纯映射 \(f_n:U_n\to\mathbb D\)，满足 \(f_n(p)=0\)、\(a_n=f_n'(p)>0\)。撇号均相对于固定坐标 \(z\) 计算。定义
\[
	r_n=\frac1{a_n},\qquad F_n=r_nf_n:U_n\longrightarrow D_{r_n},
	\qquad F_n(p)=0,\quad F_n'(p)=1,
\]
其中 \(D_r=\{w:|w|<r\}\)。\(r_n\) 是这个坐标约定下的\term[gongxingbanjing]{共形半径}[conformal radius]。

\begin{proposition}
	\label{g4:prop:0503:radii}
	对于上述圆盘穷竭，\(r_n\) 严格递增，存在极限 \(R\in(0,\infty]\)。若 \(n\ge m\)，记
	\[
		\phi_m(\zeta)=F_m^{-1}(r_m\zeta),\qquad
		t_{nm}(\zeta)=\frac{F_n(\phi_m(\zeta))}{r_n}.
	\]
	则 \(t_{nm}:\mathbb D\to\mathbb D\) 单叶，\(t_{nm}(0)=0\)，且 \(t_{nm}'(0)=r_m/r_n\)。
\end{proposition}
\begin{proof}
	\(t_{nm}=f_n\circ f_m^{-1}\) 是圆盘到自身的全纯映射。Schwarz 引理给出 \(a_n/a_m\le1\)，故 \(r_m\le r_n\)。若 \(n>m\) 而等号成立，则该映射为旋转，正导数使它为恒等。因此 \(f_n(U_m)=\mathbb D=f_n(U_n)\)，与 \(U_m\subsetneq U_n\) 及单射性矛盾。公式中的导数由链式法则得到。由于 \(r_1>0\)，递增极限严格为正。
\end{proof}

这些映照的限制未必逐点相同。它们的相容性表现为过渡映射固定原点且导数已知；下一步要把这一阶信息变成整个紧集上的控制。

\subsection{正规族极限及非退化性}
\label{g4:subsec:0503:03}

\begin{lemma}[面积定理]
	\label{g4:lem:0503:area}
	设 \(G\) 在 \(|\zeta|>1\) 单叶全纯，并在无穷远有展开
	\[
		G(\zeta)=\zeta+b_0+\sum_{k\ge1}b_k\zeta^{-k}.
	\]
	则 \(\sum_{k\ge1}k|b_k|^2\le1\)。
\end{lemma}
\begin{proof}
	固定 \(r>1\)。像曲线 \(G(re^{\mathrm{i}t})\) 为光滑 Jordan 曲线；由于无穷远处首项为 \(\zeta\)，按 \(t\) 递增取向时，其左侧为有界内部。平面 Green 公式\eqref{g4:eq:0101:green}给出该内部面积
	\[
		\frac1{2\mathrm{i}}\int_{|\zeta|=r}\overline{G(\zeta)}\,G'(\zeta)\,\mathrm{d}\zeta
		=\pi\left(r^2-\sum_{k\ge1}k|b_k|^2r^{-2k}\right)\ge0.
	\]
	Laurent 级数及导数在这个圆周上一致收敛，故逐项积分合法；不同频率的积分为零，常数项也不贡献。于是每个有限部分和满足 \(\sum_{k=1}^Nk|b_k|^2r^{-2k}\le r^2\)。先令 \(r\downarrow1\)，再令 \(N\to\infty\)，即得结论。
\end{proof}

\begin{lemma}
	\label{g4:lem:0503:distortion}
	设 \(f:\mathbb D\to\mathbb C\) 单叶全纯，\(f(0)=0\)、\(f'(0)=1\)。写 \(f(z)=z+a_2z^2+\cdots\)，则 \(|a_2|\le2\)，且对 \(|z|=r<1\)
	\[
		\frac{1-r}{(1+r)^3}\le |f'(z)|\le\frac{1+r}{(1-r)^3},
		\qquad |f(z)|\le\frac r{(1-r)^2}.
	\]
	此外 \(D_{1/4}\subset f(\mathbb D)\)。
\end{lemma}
\begin{proof}
	函数 \(f(z)/z\) 在圆盘上全纯且不为零，因而有取值为 \(1\) 的平方根。令 \(h(z)=\sqrt{f(z^2)}=z+(a_2/2)z^3+\cdots\)，其中选择 \(h'(0)=1\)。它是奇函数，且单叶：\(h(z_1)=h(z_2)\) 先给出 \(z_1^2=z_2^2\)，若两点互为相反数，则奇性迫使两者均为零。因此
	\[
		G(\zeta)=\frac1{h(1/\zeta)}=\zeta-\frac{a_2}{2\zeta}+\cdots
	\]
	满足面积定理，由此 \(|a_2|\le2\)。这里平方根的存在来自\cref{g4:prop:0102:logarithm}在单连通域上对全纯对数的构造。

	对任意 \(a\in\mathbb D\)，把 \(f\) 改成仍归一化的单叶函数
	\[
		\frac{f((\zeta+a)/(1+\overline a\zeta))-f(a)}{(1-|a|^2)f'(a)}.
	\]
	它的二次系数为 \((1-|a|^2)f''(a)/(2f'(a))-\overline a\)。刚得的界给出
	\[
		\left|\frac{f''(a)}{f'(a)}-\frac{2\overline a}{1-|a|^2}\right|
		\le\frac4{1-|a|^2}.
	\]
	沿射线 \(a=re^{\mathrm{i}\vartheta}\) 取实部，得到 \(\mathrm{d}\log|f'(re^{\mathrm{i}\vartheta})|/\mathrm{d}r\) 介于 \((2r-4)/(1-r^2)\) 和 \((2r+4)/(1-r^2)\) 之间。由 \(f'(0)=1\) 积分，便得两侧导数界。再沿同一射线积分上界，得到 \(|f(z)|\le\int_0^r(1+s)/(1-s)^3\,\mathrm{d}s=r/(1-r)^2\)。

	若 \(w\ne0\) 不在像中，则 \(wf/(w-f)\) 仍是归一化单叶函数，其二次系数为 \(a_2+1/w\)。故 \(|1/w|\le|a_2+1/w|+|a_2|\le4\)。所有遗漏值都满足 \(|w|\ge1/4\)，这就证明了最后的包含关系。
\end{proof}

最后的结论称为 Koebe 四分之一圆盘定理。\footnote{Paul Koebe（1882--1945），德国数学家，研究共形映照与黎曼曲面的统一化。}本节只需要增长上界和这个内含圆盘；它们分别控制正映射和逆映射。

\begin{lemma}
	\label{g4:lem:0503:near-identity}
	设 \(t:\mathbb D\to\mathbb D\) 全纯，\(t(0)=0\)，\(t'(0)=a\in(0,1]\)。则对 \(|z|\le s<1\)
	\[
		|t(z)-z|\le s(1-a)\frac{1+s}{1-s}.
	\]
\end{lemma}
\begin{proof}
	Schwarz 引理使 \(q(z)=t(z)/z\) 可去延拓为到闭单位圆盘的全纯映射，\(q(0)=a\)。若 \(a=1\)，最大值原理给出 \(q=1\)。否则对 \(q\) 应用 Schwarz--Pick，有 \(q(z)=(a+v)/(1+av)\)，其中 \(|v|\le|z|\)。于是 \(|q(z)-1|=(1-a)|v-1|/|1+av|\le(1-a)(1+s)/(1-s)\)，再乘 \(|z|\) 即可。
\end{proof}

\begin{lemma}
	\label{g4:lem:0503:inverse-compactness}
	在\cref{g4:prop:0503:radii}的记号下，对任意 \(0<r<R\)，存在紧集 \(K\subset X\) 及 \(n_0\)，使所有 \(n\ge n_0\) 的逆映射 \(F_n^{-1}\) 在 \(\overline D_r\) 上有定义，且像包含在 \(K\) 中。
\end{lemma}
\begin{proof}
	先设 \(R=\infty\)。选 \(m\) 使 \(r_m>4r\)。对 \(n\ge m\)，函数 \(F_n\circ\phi_m/r_m\) 在单位圆盘上单叶且导数为 \(1\)。由四分之一圆盘定理，\(D_{r_m/4}\subset F_n(U_m)\)。所以 \(H_{nm}=\phi_m^{-1}\circ F_n^{-1}\) 在 \(D_{r_m/4}\) 上定义，取值于 \(\mathbb D\)，且固定原点。Schwarz 引理给出
	\[
		|H_{nm}(w)|\le\frac{4|w|}{r_m}.
	\]
	因此 \(K=\phi_m(\overline D_{4r/r_m})\) 就满足要求。

	再设 \(R<\infty\)。选 \(0<s<1\)，使 \(r/R<s\)。当 \(m\) 足够大时，\(r_m\) 接近 \(R\)，从而对一切 \(n\ge m\)，都有 \(a=r_m/r_n\ge r_m/R\) 接近 \(1\)。\cref{g4:lem:0503:near-identity}于是保证在 \(|\zeta|=s\) 上
	\[
		|t_{nm}(\zeta)-\zeta|<s-\frac r{r_m}.
	\]
	右边可以为正，因为 \(r_m\to R\)。对任意 \(|w|\le r\)，有 \(|w/r_n|\le r/r_m\)；Rouché 定理说明 \(t_{nm}(\zeta)-w/r_n\) 在 \(D_s\) 有一个零点。故 \(F_n^{-1}(w)\in\phi_m(D_s)\)，可取紧集 \(K=\phi_m(\overline D_s)\)。
\end{proof}

\begin{theorem}
	\label{g4:thm:0503:noncompact-uniformization}
	设 \(X\) 是单连通非紧黎曼曲面，\(p\in X\)，并固定 \(p\) 附近的中心坐标。则上述 \(F_n\) 在 \(X\) 的紧集上局部一致收敛于双全纯映射
	\[
		F:X\longrightarrow D_R,
		\qquad F(p)=0,\quad F'(p)=1,
	\]
	其中 \(D_\infty=\mathbb C\)。
\end{theorem}
\begin{proof}
	固定 \(m\)。对 \(n\ge m\)，归一化单叶函数 \(F_n\circ\phi_m/r_m\) 满足\cref{g4:lem:0503:distortion}的增长上界。因此 \(F_n\) 在 \(U_m\) 的每个紧集上有与 \(n\) 无关的界。在可数张相对紧坐标图中应用\cref{g4:thm:0104:montel}并取对角子列，就得到 \(X\) 上的全纯极限 \(F\)。Cauchy 公式使导数也收敛，所以 \(F'(p)=1\)，极限不是常数。

	极限单射也可在曲面上直接验证。若两个不同点 \(x,y\) 有 \(F(x)=F(y)\)，在它们的不交小坐标圆盘中，\(F-F(x)\) 各有零点。局部一致收敛与 Rouché 定理使充分大的 \(F_n-F(x)\) 在两个圆盘中均有零点，违反 \(F_n\) 单射。故 \(F\) 单射，局部幂模型又使 \(F'\) 处处非零。若 \(R<\infty\)，由 \(|F_n|<r_n\le R\) 及最大值原理，\(F(X)\subset D_R\)；\(R=\infty\) 时这一包含自动成立。

	现在证明像没有遗漏。固定 \(0<r<R\)。逆映射紧性引理使 \(F_n^{-1}(\overline D_r)\) 全部落在同一个 \(\phi_m(\overline D_s)\) 中，其中 \(s<1\)。函数 \(H_n=\phi_m^{-1}\circ F_n^{-1}\) 在一个略大于 \(\overline D_r\) 的圆盘上有界；这是把该引理先用于稍大的半径所得。再抽子列，\(H_n\) 收敛于全纯函数 \(H\)，且在 \(\overline D_r\) 上取值于 \(\overline D_s\)。于是 \(h=\phi_m\circ H\) 定义良好。对恒等式 \(F_n(F_n^{-1}(w))=w\) 取极限，紧集上的一致收敛给出 \(F(h(w))=w\)。因此 \(D_r\subset F(X)\)。令 \(r\uparrow R\)，得到 \(F(X)=D_R\)。

	最后说明整列收敛。任意子列都可按上述办法抽到满足同样归一化的双全纯极限。有限 \(R\) 时，两个极限相差一个固定原点、导数为 \(1\) 的圆盘自同构，故相同。无限 \(R\) 时，整平面的单叶全纯函数必为仿射函数：在任意中心 \(a\) 的半径 \(L\) 圆盘上归一化，二次系数界给出 \(L|f''(a)|/(2|f'(a)|)\le2\)；令 \(L\to\infty\) 得 \(f''(a)=0\)。所以两个极限同样因值与导数归一化而相同。正规性和极限唯一性遂使整列收敛。
\end{proof}

\subsection{双曲型与抛物型的区分}
\label{g4:subsec:0503:04}

单连通非紧曲面若双全纯于 \(\mathbb D\)，称为\term[shuangquxing]{双曲型}[hyperbolic type]；若双全纯于 \(\mathbb C\)，称为\term[paowuxing]{抛物型}[parabolic type]。前一定理中的 \(R<\infty\) 与 \(R=\infty\) 分别对应这两类。二者不可能同构：如果 \(\mathbb C\to\mathbb D\) 是双全纯映射，它是有界的非常值整函数，与 Cauchy 估计推出的 Liouville 定理矛盾。\footnote{Joseph Liouville（1809--1882），法国数学家，构造了超越数，并研究 Hamilton 动力学的体积保持性质。}

\ProseParagraphBreak

共形半径的具体数值依赖中心坐标，有限或无限却不依赖。若换坐标 \(\zeta=cz+O(z^2)\)、\(c\ne0\)，则半径都乘以 \(|c|\)。换穷竭时，任一紧闭盘最终包含在另一穷竭的某个圆盘中，Schwarz 引理的半径单调性从两侧比较其上确界；故极限半径不变。换基点虽然没有这么简单的数值关系，但所得整曲面与圆盘或平面的双全纯类型已经唯一确定。

\begin{example}
	\label{g4:exm:0503:radius-models}
	在 \(\mathbb C\) 中用 \(U_n=D_n\) 穷竭、以 \(z\) 为中心坐标，则 \(F_n(z)=z\)、\(r_n=n\)，得到抛物型。在 \(\mathbb D\) 中用 \(U_n=D_{1-1/(n+1)}\)，同样有 \(F_n(z)=z\)，但 \(r_n\uparrow1\)，得到双曲型。两种情况下局部坐标都不发生变化，区别是它能扩展到的归一化半径。
\end{example}

\subsection{紧单连通情形和黎曼球面}
\label{g4:subsec:0503:05}

\begin{theorem}[单连通曲面的统一化]
	\label{g4:thm:0503:uniformization}
	每个连通、单连通的黎曼曲面恰与下列三者之一双全纯：
	\[
		\mathbb P^1,\qquad \mathbb C,\qquad\mathbb D.
	\]
\end{theorem}
\begin{proof}
	非紧情形已证。若 \(X\) 紧，则它是闭定向曲面。第一卷第6.4--6.5节的紧曲面分类及标准多边形计算给出 \(H_1(X;\mathbb Z)\cong\mathbb Z^{2g}\)；单连通使一阶同调为零，故 \(g=0\)。由\cref{g4:thm:0305:low-genus}，\(X\cong\mathbb P^1\)。具体地，Riemann--Roch 使任一点 \(p\) 的空间 \(H^0(\mathcal O([p]))\) 为二维，故存在仅有一个简单极点的非常值函数；其到球面的次数为 \(1\)，从而为双全纯同构。球面由紧性与另两者区别，平面与圆盘则由 Liouville 定理区别。
\end{proof}

这个证明的紧性分界是真实的：紧曲面上的全纯函数只能为常数，所以紧情形要构造到球面的亚纯函数；非紧情形则通过圆盘穷竭得到一个全局单叶的全纯坐标。

\subsection{万有覆盖与曲面的商表示}
\label{g4:subsec:0503:06}

一般曲面的环路未必可缩，但在万有覆盖上可以。第一卷第5.4节定理5.4.4用道路类构造万有覆盖，第5.5节命题5.5.3说明覆盖变换群自由地作用于覆盖空间，并在每个纤维上传递。黎曼曲面每点有可缩坐标圆盘，满足其中的局部道路连通与半局部单连通假设。

\begin{theorem}
	\label{g4:thm:0503:quotient}
	设 \(X\) 是连通黎曼曲面。其万有覆盖有唯一使投影局部双全纯的复结构，并且
	\[
		X\cong\Omega/\Gamma,\qquad
		\Omega\in\{\mathbb P^1,\mathbb C,\mathbb D\},
	\]
	其中 \(\Gamma\) 是 \(\Omega\) 的全纯自同构群的一个子群，其作用自由且适当不连续。这里适当不连续指每个点都有开邻域 \(V\)，满足 \(\gamma V\cap V=\varnothing\) 对所有 \(\gamma\ne1\) 成立。\(\Omega\) 的双全纯类型由 \(X\) 唯一确定。
\end{theorem}
\begin{proof}
	在均匀覆盖片上提升 \(X\) 的全纯坐标。交叠的坐标变换来自底空间，因而全纯；投影的局部逆也使这个复结构唯一。覆盖空间 Hausdorff：不同投影点用底空间的不交邻域分离，同一投影点上不同片用覆盖邻域分离。它亦第二可数。为说明这一点，紧曲面由有限多边形直接给出有限生成的基本群；非紧曲面由紧区域穷竭得到 \(\pi_1(X)\) 可数：每条环路都位于某个紧区域中，每个这样的区域基本群有限生成，故其像的可数并覆盖全部基本群。于是一个可数坐标基的每个成员只有可数张提升片，这些片的可数坐标基组成总空间的可数基。

	对这个单连通黎曼曲面应用\cref{g4:thm:0503:uniformization}，得到 \(\Omega\)。覆盖变换在提升坐标中全纯；第一卷所证的纤维传递性与局部不交性质给出 \(X=\Omega/\Gamma\)，商坐标与原坐标一致，所以同构是双全纯的。万有覆盖在指定一个提升基点后由提升唯一性唯一确定；其覆盖同构在这些坐标中全纯，因此 \(\Omega\) 的类型不变。
\end{proof}

球面型的商只有球面本身。事实上，球面的全纯自同构为一次分式变换；其不动点方程在射影直线上总有解。自由作用因此不容许任何非恒等元素。平面型的覆盖变换则都是平移：整平面的自同构由前述二次系数界为 \(z\mapsto az+b\)，而 \(a\ne1\) 时有不动点。于是平面型曲面为 \(\mathbb C/\Lambda\)，其中 \(\Lambda\) 为离散加法子群。

\ProseParagraphBreak

这样的子群只有秩 \(0,1,2\)。简要核对其格结构：若全部元素共线，离散性给出最短正向元素，带余除法说明它生成子群；若存在不共线元素，先取最短非零元素 \(\lambda_1\)，再将其他元素沿 \(\lambda_1\) 化到一个固定宽度的条带。正横向距离的下确界严格为正且取到，否则在这个条带的有界部分会出现积聚点。取取得最小横向距离的 \(\lambda_2\)，先减去其整数倍，再减去 \(\lambda_1\) 的整数倍，余项由最小性必须为零。所以 \(\Lambda=\mathbb Z\lambda_1+\mathbb Z\lambda_2\)。相应曲面是平面、\(\mathbb C^\times\) 和复环面；秩一时指数 \(\exp(2\pi\mathrm{i}z/\lambda_1)\) 给出第二个识别。只有秩二的商紧。

\ProseParagraphBreak

因此紧曲面的类型与亏格相配：亏格零为球面，亏格一由\cref{g4:cor:0403:abel-embedding}为复环面，亏格至少二的万有覆盖只能是圆盘。一般非紧曲面也按其万有覆盖称为双曲型或抛物型；例如穿孔圆盘是双曲型，而 \(\mathbb C^\times\) 是抛物型。它们都具有无限循环基本群，说明基本群本身不足以决定共形类型。下一节将给双曲型曲面一个由复结构唯一确定的完备度量。
% 纲要条目排为编号小节标题，并进入目录与书签；说明段仍保留于纲要注释中。

% !TeX root = ../../Book4.tex
% 纲要标题：### 5.4 共形度量与双曲收缩
% 纲要位置：几何学计划纲要.md，第 3663 行。
% 正文状态：初稿；尚未正式审核。

% 纲要内容开始
% 
% * 曲率的共形变换
% * Schwarz–Pick 定理与 Ahlfors–Schwarz 引理
% * Poincaré 度量及完备性
% * 尖点、环域与漏斗
% 
% 纲要内容结束

\section{共形度量与双曲收缩}
\label{g4:sec:0504}

统一化给出圆盘、平面和球面三种全纯模型，但坐标分类还没有说明如何量长度。圆盘上的 Schwarz--Pick 不等式已经指出一个特别的长度密度：全纯映射对于它总是收缩的。本节从曲率计算解释这一选择，再把圆盘度量下降到双曲型曲面。尖点和漏斗则显示，同样的局部负曲率可以产生不同的整体端部。

\subsection{曲率的共形变换}
\label{g4:subsec:0504:01}

在全纯坐标 \(z=x+\mathrm{i}y\) 中，共形度量写成 \(g=\rho^2|\mathrm{d}z|^2\)，其中 \(\rho>0\) 光滑。记 \(\Delta_{\mathrm e}=\partial_x^2+\partial_y^2\)，本书的非负 Laplace 算子则是 \(\Delta_g=-\rho^{-2}\Delta_{\mathrm e}\)。\(\rho\) 是长度的缩放因子，面积的缩放因子为 \(\rho^2\)。若另一个全纯坐标为 \(w=w(z)\)，长度密度满足 \(\rho_z=\rho_w|w'|\)，因此这些局部表达拼成同一个度量。

\begin{proposition}
	\label{g4:prop:0504:conformal-curvature}
	设 \(X\) 为黎曼曲面，\(g\) 为光滑正定共形度量，\(\varphi\in C^\infty(X,\mathbb R)\)。采用标准球面曲率为正及非负 Laplace 的约定，有
	\begin{equation}
		\label{g4:eq:0504:curvature-change}
		K_g=-\rho^{-2}\Delta_{\mathrm e}\log\rho,
		\qquad K_{e^{2\varphi}g}=e^{-2\varphi}(K_g+\Delta_g\varphi).
	\end{equation}
\end{proposition}
\begin{proof}
	第一卷第8.3节例8.3.10由度量的联络计算了第一个公式。为固定这里的符号，令 \(u=\log\rho\)，恢复计算中所需的系数：
	\[
		\begin{array}{lll}
			\Gamma^1_{11}=u_x,&\Gamma^1_{12}=u_y,&\Gamma^1_{22}=-u_x,\\
			\Gamma^2_{11}=-u_y,&\Gamma^2_{12}=u_x,&\Gamma^2_{22}=u_y.
		\end{array}
	\]
	在 \(R(\partial_x,\partial_y)\partial_y\) 的 \(\partial_x\) 分量中，一阶导数乘积相消，剩下 \(-u_{xx}-u_{yy}\)。与 \(\partial_x\) 作内积后除以坐标平行四边形面积的平方 \(e^{4u}\)，即得 \(K_g=-e^{-2u}\Delta_{\mathrm e}u\)。将 \(u\) 换成 \(u+\varphi\)，再用 \(\Delta_g\varphi=-e^{-2u}\Delta_{\mathrm e}\varphi\)，便得到第二式。换全纯坐标时，\(\log|w'|\) 调和，所以公式与坐标变换相容。
\end{proof}

\begin{example}
	\label{g4:exm:0504:three-curvatures}
	平面上的密度 \(\rho=1\) 给出 \(K=0\)。球极坐标中的密度 \(\rho=2/(1+|z|^2)\) 给出 \(K=1\)，并在无穷远换坐标 \(w=1/z\) 后保持同一形式。单位圆盘上的密度 \(\lambda_{\mathbb D}=2/(1-|z|^2)\) 则满足
	\[
		\Delta_{\mathrm e}\log\lambda_{\mathbb D}
		=\frac4{(1-|z|^2)^2}=\lambda_{\mathbb D}^2,
	\]
	所以曲率为 \(-1\)。由 \(C(z)=(z-\mathrm{i})/(z+\mathrm{i})\) 将上半平面映到圆盘，直接计算
	\[
		\lambda_{\mathbb D}(C(z))|C'(z)|=\frac1{\operatorname{Im}z}.
	\]
	因而上半平面的对应度量是 \(|\mathrm{d}z|^2/y^2\)，也有曲率 \(-1\)。
\end{example}

若在固定共形类内寻找曲率 \(-1\) 的度量，上式把问题写成 \(\Delta_g\varphi+K_g=-e^{2\varphi}\)。统一化在双曲型曲面上给出了一个全局解的几何来源。下面的比较则说明，负曲率条件本身已经限制了长度密度能够有多大。

\subsection{Schwarz–Pick 定理与 Ahlfors–Schwarz 引理}
\label{g4:subsec:0504:02}

\cref{g4:thm:0103:schwarz-pick}已从 Schwarz 引理证明，对全纯映射 \(f:\mathbb D\to\mathbb D\)，有
\[
	\lambda_{\mathbb D}(f(z))|f'(z)|\le\lambda_{\mathbb D}(z).
\]
左边是目标双曲度量拉回后的长度密度；分歧点处它可以为零。要把圆盘目标换成任意具有负曲率上界的曲面，需要一个允许这种零点的比较结论，即 Ahlfors\footnote{Lars Valerian Ahlfors（1907—1996），芬兰数学家，研究复分析与黎曼曲面。}--Schwarz 引理。

\begin{theorem}[Ahlfors--Schwarz 引理]
	\label{g4:thm:0504:ahlfors-schwarz}
	设 \(\mu:\mathbb D\to[0,\infty)\) 连续，在开集 \(\{\mu>0\}\) 上属于 \(C^2\)，并在该开集上满足
	\[
		\Delta_{\mathrm e}\log\mu\ge\mu^2.
	\]
	则对每个 \(z\in\mathbb D\)，有
	\[
		\mu(z)\le\frac2{1-|z|^2}.
	\]
	特别地，若 \(Y\) 为具有光滑正定共形度量 \(h\) 的黎曼曲面，\(K_h\le-1\)，则任意全纯映射 \(f:\mathbb D\to Y\) 满足 \(f^*h\le g_{\mathbb D}\)。
\end{theorem}
\begin{proof}
	先在较小圆盘 \(|z|<R<1\) 上比较，取
	\[
		\lambda_R(z)=\frac{2R}{R^2-|z|^2}.
	\]
	直接求导有 \(\Delta_{\mathrm e}\log\lambda_R=\lambda_R^2\)，且 \(\lambda_R\) 在圆周 \(|z|=R\) 趋于无穷大。\(\mu\) 在闭圆盘 \(|z|\le R\) 有界，所以 \(\mu/\lambda_R\) 连续延拓到边界，边界值为零。

	如果这个比值在某处大于 \(1\)，其最大值就在内部某点 \(z_0\) 取得，且 \(\mu(z_0)>0\)。故 \(\log(\mu/\lambda_R)\) 在那里可以二次求导，并满足
	\[
		0\ge\Delta_{\mathrm e}\log\frac\mu{\lambda_R}(z_0)
		\ge\mu(z_0)^2-\lambda_R(z_0)^2>0,
	\]
	矛盾。因此 \(\mu\le\lambda_R\)；固定 \(z\) 并让 \(R\uparrow1\)，即得所求界。允许 \(\mu\) 消失并不妨碍证明，因为假设的正最大值一定出现在 \(\mu>0\) 的区域。

	对最后的断言，常值映射显然满足结论。非常值全纯映射在目标坐标中给出密度 \(\mu(z)=\rho_Y(f(z))|f'(z)|\)，它连续，且在 \(f'\ne0\) 处光滑。由于 \(\log|f'|\) 在这里调和，链式法则给出
	\[
		\Delta_{\mathrm e}\log\mu
		=(\Delta_{\mathrm e,w}\log\rho_Y)(f(z))|f'(z)|^2
		=-K_h(f(z))\mu^2\ge\mu^2.
	\]
	这些表达在目标坐标交叠处相容。应用前半部分，得到 \(\mu\le\lambda_{\mathbb D}\)，平方后就是度量比较。
\end{proof}

目标度量的完备性没有出现在这个证明中；比较利用的是曲率不等式和源圆盘的边界发散。若目标曲率满足 \(K_h\le-\kappa\)、\(\kappa>0\)，将 \(h\) 乘以 \(\kappa\) 后曲率至多为 \(-1\)，所以结论相应变为 \(f^*h\le\kappa^{-1}g_{\mathbb D}\)。这个尺度因子在比较不同曲率归一化时不可省略。

\begin{corollary}
	\label{g4:cor:0504:entire-negative-curvature}
	设 \(Y\) 为具有光滑正定共形度量 \(h\) 的黎曼曲面，且 \(K_h\le-\kappa<0\)。则每个全纯映射 \(F:\mathbb C\to Y\) 都是常值。
\end{corollary}
\begin{proof}
	固定 \(z_0\in\mathbb C\)，将上一定理用于 \(\zeta\mapsto F(z_0+R\zeta)\)。在 \(\zeta=0\) 得到
	\[
		R\rho_Y(F(z_0))|F'(z_0)|\le\frac2{\sqrt\kappa}.
	\]
	令 \(R\to\infty\)，知 \(F'(z_0)=0\)。每点如此，连通性使 \(F\) 为常值。
\end{proof}

\subsection{Poincaré 度量及完备性}
\label{g4:subsec:0504:03}

称黎曼曲面为双曲型，是指其万有覆盖双全纯于 \(\mathbb D\)。由\cref{g4:thm:0503:quotient}，取全纯覆盖 \(\pi:\mathbb D\to X\)，各覆盖变换都是圆盘自同构。对自同构及其逆分别应用 Schwarz--Pick，得到长度密度的变换等式，因此圆盘度量能够沿 \(\pi\) 下降。

\begin{definition}
	\label{g4:def:0504:poincare}
	设 \(X\) 为双曲型连通黎曼曲面。唯一满足
	\[
		\pi^*g_X=g_{\mathbb D}
	\]
	的光滑共形度量 \(g_X\) 称为它的\term[Poincareduliang]{Poincaré 度量}[Poincaré metric]\footnote{Jules Henri Poincaré（1854—1912），法国数学家，研究微分方程、拓扑与天体力学。}。其长度距离记为 \(d_X\)。
\end{definition}

这里的唯一性是相对于覆盖拉回等式而言。局部取 \(\pi\) 的逆分支便能定义度量，换逆分支只相差圆盘自同构，故与选择无关。换用另一张万有覆盖时，两张覆盖之间也有圆盘自同构，所得度量仍相同。局部拉回保持曲率，所以 \(K_{g_X}=-1\)。

\begin{theorem}
	\label{g4:thm:0504:poincare-complete}
	每个双曲型连通黎曼曲面的 Poincaré 长度距离完备。任意全纯映射 \(F:X\to Y\)，其中 \(X,Y\) 均为双曲型连通黎曼曲面，都满足
	\[
		F^*g_Y\le g_X,
		\qquad d_Y(F(p),F(q))\le d_X(p,q).
	\]
	若第一式在某点的一个非零切向量上取等号，则 \(F\) 是全纯覆盖；反之每个全纯覆盖都是上述度量的局部等距。
\end{theorem}
\begin{proof}
	先证明 \(\mathbb D\) 完备。\cref{g4:prop:0103:hyperbolic-distance-formula}给出
	\[
		d_{\mathbb D}(0,z)=\log\frac{1+|z|}{1-|z|}.
	\]
	一个双曲 Cauchy 列距原点的距离有界，因而留在某个闭子圆盘 \(|z|\le r<1\) 内。双曲距离又至少为两倍 Euclid 距离，所以该列有 Euclid 极限 \(z\)，仍在这个闭子圆盘中。密度在稍大的闭子圆盘上有界，连接 \(z_n,z\) 的直线段便说明它也按双曲距离收敛。

	再设 \((p_n)\) 为 \(X\) 中的 Cauchy 列。取子列，使相邻项可由长度小于 \(2^{-k}\) 的分片光滑曲线 \(\gamma_k\) 连接；这是长度距离的定义与 Cauchy 条件共同保证的。在 \(\mathbb D\) 中选定第一项的原像，由第一卷第5.3节定理5.3.5的道路提升依次提升 \(\gamma_k\)。由于 \(\pi\) 是局部等距，提升曲线长度不变，其端点列 \((z_k)\) 满足
	\[
		d_{\mathbb D}(z_k,z_j)\le\sum_{m=k}^{j-1}2^{-m}\quad(j>k).
	\]
	于是 \(z_k\) 收敛于某个 \(z\in\mathbb D\)，投影后原子列收敛于 \(\pi(z)\)。原列 Cauchy，故整列也收敛。这里长度距离给出原流形拓扑：在相对紧坐标小片上，光滑正密度与 Euclid 密度双边可比，而离开较小片的曲线有一个固定的正长度下界。

	对全纯映射 \(F\)，第一卷第5.3节定理5.3.10的提升判据适用于道路连通、局部道路连通且单连通的圆盘，基本群像包含条件自动成立。因此组合 \(F\circ\pi_X:\mathbb D\to Y\) 提升为全纯映射 \(f:\mathbb D\to\mathbb D\)，满足 \(\pi_Yf=F\pi_X\)。提升的全纯性由覆盖的局部全纯逆得到。对 \(f\) 用 Schwarz--Pick，再用两边覆盖的局部等距性质，得到 \(F^*g_Y\le g_X\)。沿连接 \(p,q\) 的任意曲线积分并取下确界，便得距离不增加。

	若某点某非零方向取等号，对应的 \(f\) 在一点满足 Schwarz--Pick 的导数等号，故 \(f\) 为圆盘自同构。把 \(X\) 的覆盖变换群经 \(f\) 共轭后，\(F\) 就是某个子群商到 \(Y\) 的群商的自然映射，因此为全纯覆盖。具体地，覆盖的局部均匀覆盖邻域在中间子群商中仍分解成互不相交的同构副本，这正给出覆盖条件。反过来，若 \(F\) 为全纯覆盖，\(F\pi_X\) 也是 \(Y\) 的万有覆盖，所以它与 \(\pi_Y\) 相差圆盘自同构，局部等距性随即成立。
\end{proof}

对于一般曲面，局部等号只推出覆盖，未必推出双射。例如 \(z\mapsto z^n\) 是穿孔圆盘的 \(n\) 重全纯覆盖，\(n\ge2\) 时并不单射。相应度量将在下一小节显式写出。

\begin{proposition}
	\label{g4:prop:0504:poincare-maximal}
	设 \(X\) 为双曲型连通黎曼曲面，\(h\) 是其任一光滑正定共形度量，且 \(K_h\le-1\)。则 \(h\le g_X\)。
\end{proposition}
\begin{proof}
	把 \(h\) 拉回到万有覆盖圆盘，应用\cref{g4:thm:0504:ahlfors-schwarz}，得到 \(\pi^*h\le g_{\mathbb D}=\pi^*g_X\)。覆盖在每点为局部微分同胚，故比较下降为 \(h\le g_X\)。
\end{proof}

\subsection{尖点、环域与漏斗}
\label{g4:subsec:0504:04}

第一卷第10.2节例10.2.10和命题10.2.11已计算抛物平移与双曲缩放产生的尖点、漏斗。本节保留同样的曲率和长度归一化，并补出它们的全纯坐标。指数映射使无限条带变成环域，因而也把端部的度量行为转化为一个可计算的共形参数。

\begin{example}
	\label{g4:exm:0504:cusp}
	在上半平面 \(\mathbb H\) 上令 \(T(z)=z+1\)。函数 \(q=e^{2\pi\mathrm{i}z}\) 给出
	\[
		\mathbb H/\langle T\rangle\cong\{0<|q|<1\}.
	\]
	指数的相等恰表示参数相差整数，且任意 \(0<|q|<1\) 都有虚部为正的对数原像，所以这是双全纯同构。由 \(|q|=e^{-2\pi y}\)、\(|\mathrm{d}q|=2\pi|q||\mathrm{d}z|\)，得到穿孔圆盘的 Poincaré 度量
	\begin{equation}
		\label{g4:eq:0504:punctured-disk-metric}
		\mathrm{d}s^2=\frac{|\mathrm{d}q|^2}{|q|^2\bigl(\log(1/|q|)\bigr)^2}.
	\end{equation}
	固定 \(Y>0\)，商中 \(y\ge Y\) 的部分称为一个\term[jiandian]{尖点}[cusp]端。写 \(r=\log(y/Y)\)，则
	\[
		\mathrm{d}s^2=\mathrm{d}r^2+Y^{-2}e^{-2r}\,\mathrm{d}x^2,
		\qquad r\ge0,\quad x\in\mathbb R/\mathbb Z.
	\]
	截圈长为 \(Y^{-1}e^{-r}\)，端部面积为 \(1/Y\)，而走向 \(r\to\infty\) 的曲线长度至少为 \(r\) 的总变化，所以尖端在无限距离处。整个穿孔圆盘的另一端 \(|q|\to1\) 也在无限距离处，但其面积无限。
\end{example}

\begin{proposition}
	\label{g4:prop:0504:hyperbolic-cylinder}
	设 \(\ell>0\)。商 \(\mathbb H/\langle z\mapsto e^\ell z\rangle\) 是完备双曲圆柱，双全纯于环域
	\[
		A_\ell=\{e^{-2\pi^2/\ell}<|q|<1\}.
	\]
	它具有坐标 \(r\in\mathbb R\)、\(t\in\mathbb R/\ell\mathbb Z\)，使
	\begin{equation}
		\label{g4:eq:0504:cylinder-metric}
		\mathrm{d}s^2=\mathrm{d}r^2+\cosh^2r\,\mathrm{d}t^2.
	\end{equation}
	其中 \(r=0\) 是长为 \(\ell\) 的闭测地线，每个半圆柱 \(r\ge0\) 或 \(r\le0\) 为一个\term[loudou]{漏斗}[funnel]。
\end{proposition}
\begin{proof}
	在上半平面取单值对数 \(w=\log z=t+\mathrm{i}\theta\)，其中 \(0<\theta<\pi\)。缩放对应 \(w\mapsto w+\ell\)，而 \(q=e^{2\pi\mathrm{i}w/\ell}\) 恰把条带模此平移识别为所示环域。两个生成元均无内点不动点；对数坐标中的平移还直接说明作用局部有限，所以它们确实给出全纯覆盖的商。

	度量变成
	\[
		\mathrm{d}s^2=\frac{\mathrm{d}t^2+\mathrm{d}\theta^2}{\sin^2\theta}.
	\]
	令 \(r=\operatorname{arsinh}(\cot\theta)\)，则 \(\sin\theta=\operatorname{sech}r\)、\(\mathrm{d}\theta=-\operatorname{sech}r\,\mathrm{d}r\)，代入即得式\eqref{g4:eq:0504:cylinder-metric}。这也恢复了第一卷中的参数式 \(x=e^t\tanh r\)、\(y=e^t\operatorname{sech}r\)。曲线 \(r=0\) 的测地方程由 \(\partial_r\cosh^2r|_{r=0}=0\) 直接满足，其长度为 \(t\) 的周期 \(\ell\)。任何绕圆柱一周的曲线都有长度至少 \(\int|\mathrm{d}t|\ge\ell\)，所以这条核心圈实现该自由同伦类的最小长度。

	到核心圈的距离为 \(|r|\)：任一连接曲线的长度至少为 \(r\) 的变化，固定 \(t\) 的径向段达到下界。漏斗截圈的长度为 \(\ell\cosh r\)，而半圆柱 \(0\le r\le R\) 的面积为 \(\ell\sinh R\)，随 \(R\to\infty\) 发散。整个商的完备性已由\cref{g4:thm:0504:poincare-complete}保证，也可直接从 \(|r|\) 的距离下界与紧的有界 \(r\) 区域得到。
\end{proof}

\begin{corollary}
	\label{g4:cor:0504:annulus-metric}
	设 \(0<a<b<\infty\)，\(L=\log(b/a)\)。环域 \(A(a,b)=\{a<|z|<b\}\) 的 Poincaré 长度密度为
	\begin{equation}
		\label{g4:eq:0504:annulus-density}
		\lambda_{A(a,b)}(z)
		=\frac{\pi}{L|z|\sin\bigl(\pi\log(|z|/a)/L\bigr)}.
	\end{equation}
	中间圆周 \(|z|=\sqrt{ab}\) 是长为 \(2\pi^2/L\) 的闭测地线；环域的两条边界均在无限距离处。
\end{corollary}
\begin{proof}
	在对数覆盖上取 \(w=s+\mathrm{i}t=\log(z/a)\)，其中 \(0<s<L\)、\(t\) 以 \(2\pi\) 为周期。映射 \(w\mapsto\exp(\pi\mathrm{i}w/L)\) 把这一条带双全纯地送到上半平面。拉回密度 \(1/y\)，得到 \((\pi/L)/\sin(\pi s/L)\)；再用 \(|\mathrm{d}w|=|\mathrm{d}z|/|z|\)，便得公式。平移 \(t\mapsto t+2\pi\) 在上半平面中对应缩放因子 \(e^{-2\pi^2/L}\)，所生成的群与正指数缩放生成的群相同。于是上一命题给出核心圈的测地性和长度 \(\ell=2\pi^2/L\)。靠近任一边界，密度与边界法向距离的倒数同阶，其径向积分发散；或者直接应用整个圆柱的完备性。
\end{proof}

尖点和漏斗都由圆盘型覆盖产生，而其共形边界不同：尖点对应一个删去的点，有限模环域的两个漏斗端则对应两个边界圆。环域中的 Euclid 圆周半径不能直接衡量这一区别。下一节以曲线族的最短长度和面积之比，提取不依赖所选共形度量的数值。
% 纲要条目排为编号小节标题，并进入目录与书签；说明段仍保留于纲要注释中。

% !TeX root = ../../Book4.tex
% 纲要标题：### 5.5 极值长度和模量
% 纲要位置：几何学计划纲要.md，第 3670 行。
% 正文状态：初稿；尚未正式审核。

% 纲要内容开始
% 
% * 曲线族的极值长度
% * 共形不变性与单调性
% * 环域模量的计算
% * 退化曲面的长颈模型
% 
% ---
% 
% 纲要内容结束

\section{极值长度和模量}
\label{g4:sec:0505}

同一共形类中的度量可以任意乘上正函数，单条曲线的长度因而不是共形不变量。不过，把一族曲线中最短的长度平方除以总面积，会消去常数缩放；再对全部共形密度取上确界，就得到只依赖复结构的量。环域上这个量可以完整计算，所得模量把“长而细的颈部”变成明确的数值条件。

\subsection{曲线族的极值长度}
\label{g4:subsec:0505:01}

设 \(X\) 为黎曼曲面，先选一个辅助光滑共形度量 \(g_0\)。本节的曲线取局部可求长的连续参数曲线，允许开参数区间的端点趋于 \(X\) 的边界或端部。长度用各紧参数子区间上的线积分之上确界定义，因而也允许无穷长。重新单调参数化不改变积分。

\begin{definition}
	\label{g4:def:0505:extremal-length}
	设 \(\Gamma\) 为黎曼曲面 \(X\) 上的一族上述曲线。对非负 Borel 函数 \(\rho\)，定义
	\[
		A_\rho=\int_X\rho^2\,\mathrm{d}A_0,
		\qquad L_\rho(\Gamma)=\inf_{\gamma\in\Gamma}\int_\gamma\rho\,\mathrm{d}s_0.
	\]
	只取满足 \(0<A_\rho<\infty\) 的密度作为可容许密度。曲线族的\term[jizhichangdu]{极值长度}[extremal length]为
	\begin{equation}
		\label{g4:eq:0505:extremal-length}
		\operatorname{EL}_X(\Gamma)
		=\sup_{0<A_\rho<\infty}\frac{L_\rho(\Gamma)^2}{A_\rho}
		\in[0,\infty].
	\end{equation}
	密度可在面积零集上取 \(+\infty\)；积分使用非负函数的通常约定。若 \(\Gamma=\varnothing\)，置 \(L_\rho(\Gamma)=+\infty\)，故极值长度为 \(+\infty\)。
\end{definition}

可容许密度总是存在。取相对紧穷竭 \(U_n\)，置 \(U_0=\varnothing\)，在各个互不相交的集合 \(U_n\setminus U_{n-1}\) 上取正常数 \(c_n\)，使 \(c_n^2\operatorname{Area}_{g_0}(U_n)\le2^{-n}\)，即可得到一个处处为正且平方可积的 Borel 密度。度量乘以常数 \(c^2\) 时，长度乘以 \(c\)、面积乘以 \(c^2\)，所以式\eqref{g4:eq:0505:extremal-length}的商不变。定义允许密度在部分区域为零，这使我们能够只给待考察的颈部赋予长度，而不把远处面积计入成本。

\ProseParagraphBreak

采用 Borel 密度还需留意面积与曲线积分的区别：在面积零集上改变密度，可能改变位于该集合内的某条曲线的长度。这里取点态密度及其线积分，不把密度仅视为 \(L^2\) 等价类。下面的 Fubini 与 Cauchy--Schwarz 估计只要求几乎处处的横截曲线满足积分界，恰好适用于这一范围。

\begin{example}
	\label{g4:exm:0505:rectangle}
	在矩形 \(Q=(0,a)\times(0,b)\) 中，令 \(\Gamma\) 为连接两条竖边的曲线族。则
	\[
		\operatorname{EL}_Q(\Gamma)=\frac ab.
	\]
	确实，对任一可容许 Euclid 密度 \(\rho\)，水平线段均在该族中，因此对几乎处处的 \(y\)，
	\[
		L_\rho(\Gamma)^2\le
		\left(\int_0^a\rho(x,y)\,\mathrm{d}x\right)^2
		\le a\int_0^a\rho(x,y)^2\,\mathrm{d}x.
	\]
	对 \(y\) 积分，得到 \(bL_\rho(\Gamma)^2\le aA_\rho\)。取常密度 \(\rho=1\) 时，每条连接曲线的长度至少为横坐标总变化 \(a\)，水平线段达到下界，而面积为 \(ab\)。上下界相合。交换横竖方向，连接横边的曲线族极值长度为 \(b/a\)。
\end{example}

这个计算显示极值过程的两面：选一族明确的测试曲线，可以给任何密度一个上界；找到达到上界的密度，才完成极值长度的计算。

\subsection{共形不变性与单调性}
\label{g4:subsec:0505:02}

\begin{proposition}
	\label{g4:prop:0505:conformal-invariance}
	极值长度不依赖辅助共形度量的选择。若 \(F:X\to Y\) 是黎曼曲面之间的双全纯映射，则
	\[
		\operatorname{EL}_Y(F\Gamma)=\operatorname{EL}_X(\Gamma).
	\]
\end{proposition}
\begin{proof}
	若换为 \(g_1=\sigma^2g_0\)、\(\sigma>0\) 光滑，则密度 \(\rho_1\) 对应旧度量下的 \(\rho_0=\sigma\rho_1\)。两者的曲线长度和面积分别相同，这个对应又是双射，所以极值相同。

	对双全纯 \(F\)，在局部坐标中，将目标密度 \(\eta\) 拉回为 \(\rho(z)=\eta(F(z))|F'(z)|\)。曲线积分的换元给出长度相同；实 Jacobian 为 \(|F'|^2\)，面积换元给出 \(A_\rho=A_\eta\)。逆映射使可容许密度一一对应，故对相等的商取上确界得到结论。一般辅助度量由前半部分归到同样的局部计算。
\end{proof}

\begin{proposition}
	\label{g4:prop:0505:monotonicity}
	设 \(\Gamma_1,\Gamma_2\) 为同一黎曼曲面上的曲线族。
	\begin{enumerate}
		\item 若 \(\Gamma_1\subseteq\Gamma_2\)，则 \(\operatorname{EL}(\Gamma_1)\ge\operatorname{EL}(\Gamma_2)\)。
		\item 若 \(\Gamma_1\) 中每条曲线都含有 \(\Gamma_2\) 中某条曲线作为子弧，则仍有 \(\operatorname{EL}(\Gamma_1)\ge\operatorname{EL}(\Gamma_2)\)。
		\item 若族 \(\Gamma\) 的所有曲线都位于开子域 \(U\subseteq X\)，则 \(\operatorname{EL}_U(\Gamma)=\operatorname{EL}_X(\Gamma)\)。
	\end{enumerate}
\end{proposition}
\begin{proof}
	前两项都由对每个密度成立的 \(L_\rho(\Gamma_1)\ge L_\rho(\Gamma_2)\) 得到。第一项因下确界在较小集合上取值较大；第二项因非负长度在取子弧时不增加。

	第三项中，把 \(U\) 上的密度在外部补零，长度和面积均不变，所以 \(\operatorname{EL}_X\ge\operatorname{EL}_U\)。反过来，将 \(X\) 上的密度限制到 \(U\)，曲线长度不变而面积不增加，故当限制面积为正时，相应商不减。若限制面积为零而长度为正，可在限制密度上加 \(\varepsilon\) 倍一个正的有限面积密度；面积趋于零，长度仍有原来的正下界，于是 \(\operatorname{EL}_U=\infty\)，比较仍成立。若长度也是零，原商为零，不影响上确界。因此得到相等。
\end{proof}

曲线族增大使极值长度减小，因为较多的候选使“最短一条”更难拉长。比较嵌套区域时，应先辨认具体哪一族包含哪一族，或每条曲线必须穿过哪些子弧；单说“区域变大”并不能确定不等式方向。

\begin{proposition}
	\label{g4:prop:0505:series-parallel}
	设 \(U_1,\ldots,U_k\) 是黎曼曲面 \(X\) 中互不相交的开集，\(\Gamma_i\) 是完全位于 \(U_i\) 的曲线族，并设 \(e_i=\operatorname{EL}_{U_i}(\Gamma_i)\in(0,\infty)\)。
	\begin{enumerate}
		\item 若曲线族 \(\Gamma\) 的每条曲线都含有各族 \(\Gamma_i\) 中的一条子弧，则
		\[
			\operatorname{EL}_X(\Gamma)\ge\sum_{i=1}^k e_i.
		\]
		\item 对可任选其中一族的并集 \(\Gamma=\bigcup_i\Gamma_i\)，有
		\[
			\frac1{\operatorname{EL}_X(\Gamma)}=\sum_{i=1}^k\frac1{e_i}.
		\]
	\end{enumerate}
\end{proposition}
\begin{proof}
	对第一项，在每个 \(U_i\) 选择可容许密度 \(\rho_i\)，使其长度下确界 \(L_i>0\)、面积为 \(A_i\)，且 \(L_i^2/A_i\) 任意接近 \(e_i\)。将 \((L_i/A_i)\rho_i\) 定义在 \(U_i\)，其余处取零，得到一个全局密度。任一 \(\gamma\in\Gamma\) 至少贡献 \(\sum_i L_i^2/A_i\) 的长度；各子弧在互不相交区域中，长度可以相加。总面积也恰为 \(\sum_i L_i^2/A_i\)，故所得比值至少为这个和。让各个比值趋于 \(e_i\)，便得串联不等式。

	第二项可以只在并集 \(\bigcup U_i\) 上讨论。任一全局密度在各部分的面积记为 \(A_i\)，长度下确界记为 \(L_i\)，而整个并集族的下确界为 \(L=\min_iL_i\)。定义给出 \(L_i^2\le e_iA_i\)；当 \(A_i=0\) 时，由上一命题证明中的小正密度扰动及 \(e_i<\infty\) 可知 \(L_i=0\)，所以这个不等式也成立。因此
	\[
		\sum_iA_i\ge\sum_i\frac{L^2}{e_i},
	\]
	得到 \(\operatorname{EL}(\Gamma)\le(\sum_i e_i^{-1})^{-1}\)。为得反向，仍选取近于极值的 \(\rho_i\)，并各自缩放到 \(L_i=1\)。此时面积任意接近 \(1/e_i\)。拼合密度后 \(L=1\)，面积趋于 \(\sum_i1/e_i\)，从而达到所求下界。
\end{proof}

第一项描述每条曲线必须依次穿越多个区域时的长度障碍；第二项描述可以任选通道时的结果。这两个计算仅在题设明确的分离区域和曲线条件下使用。

\subsection{环域模量的计算}
\label{g4:subsec:0505:03}

设 \(A(a,b)=\{a<|z|<b\}\)，\(0<a<b<\infty\)。记 \(\Gamma_{\mathrm{join}}\) 为从一个边界圆趋向另一个边界圆的曲线族，\(\Gamma_{\mathrm{core}}\) 为绕原点次数非零的闭曲线族。后者包括简单核心圆及其多重绕行；为求最短长度，一次绕行已经给出测试曲线。对数坐标 \(w=\log(z/a)=s+\mathrm{i}t\) 把环域变成圆柱
\[
	0<s<L=\log(b/a),\qquad t\in\mathbb R/2\pi\mathbb Z.
\]
\cref{g4:fig:0505:annulus-cylinder}显示两族曲线在这个坐标中的方向。

% !TeX root = ../Book4.tex
\begin{figure}[htbp]
	\centering
	\begin{tikzpicture}[font=\small,scale=0.9]
		\begin{scope}[xshift=-3.5cm]
			\fill[GeometryBlueLight,even odd rule] (0,0) circle (1.4) (0,0) circle (0.55);
			\draw[GeometryGray] (0,0) circle (1.4);
			\draw[GeometryGray] (0,0) circle (0.55);
			\draw[GeometryBlue,thick,->] (0.55,0)--(1.4,0);
			\draw[GeometryBlue,dashed,thick,->] (45:0.95) arc (45:365:0.95);
			\node[below] at (0,-1.65) {环域 \(a<|z|<b\)};
			\node[above] at (0,1.65) {径向连接与核心闭路};
		\end{scope}
		\draw[->] (-1.6,0)--(0.1,0) node[midway,above] {\(w=\log(z/a)\)};
		\begin{scope}[xshift=0.7cm,yshift=-1.2cm]
			\fill[GeometryBlueLight] (0,0) rectangle (3.6,2.4);
			\draw[GeometryGray] (0,0) rectangle (3.6,2.4);
			\draw[GeometryBlue,thick,->] (0,1.05)--(3.6,1.05);
			\draw[GeometryBlue,dashed,thick,->] (2,0)--(2,2.4);
			\draw[->] (0.8,0)--(1.2,0);
			\draw[->] (0.8,2.4)--(1.2,2.4);
			\node[below] at (0,0) {\(0\)};
			\node[below] at (3.6,0) {\(L=\log(b/a)\)};
			\node[left] at (0,2.4) {\(2\pi\)};
			\node[above] at (1.8,2.75) {上下两边按平移识别};
			\node[below] at (1.8,-0.6) {连接族沿 \(s\)，核心族沿 \(t\)};
		\end{scope}
	\end{tikzpicture}
	\caption[环域的对数圆柱与两族曲线]{对数坐标使密度 \(|\mathrm{d}z|/|z|\) 成为平坦长度 \(|\mathrm{d}w|\)。实线表示连接两个边界的曲线，虚线表示核心闭路；右图的上下边识别后，虚线的两端接成一圈。图示尺寸只表达坐标方向，模量由 \(L/(2\pi)\) 给出。}
	\label{g4:fig:0505:annulus-cylinder}
\end{figure}


\begin{theorem}
	\label{g4:thm:0505:annulus-modulus}
	对有限圆环域 \(A(a,b)\)，上述两族曲线满足
	\begin{equation}
		\label{g4:eq:0505:annulus-extremal-length}
		\operatorname{EL}(\Gamma_{\mathrm{join}})
		=\frac1{2\pi}\log\frac ba,
		\qquad
		\operatorname{EL}(\Gamma_{\mathrm{core}})
		=\frac{2\pi}{\log(b/a)}.
	\end{equation}
	两者都由长度密度 \(1/|z|\) 达到极值。定义环域的\term[moliang]{模量}[modulus]为
	\[
		\operatorname{Mod}A(a,b)\coloneqq\frac1{2\pi}\log\frac ba.
	\]
\end{theorem}
\begin{proof}
	在对数圆柱上用平坦辅助度量 \(|\mathrm{d}w|^2\)。对任一可容许密度 \(\rho(s,t)\)，固定 \(t\) 的水平弧属于连接族，所以 Cauchy--Schwarz 与 Fubini 给出
	\[
		2\pi L_\rho(\Gamma_{\mathrm{join}})^2
		\le L\int_0^{2\pi}\int_0^L\rho(s,t)^2\,\mathrm{d}s\,\mathrm{d}t
		=LA_\rho.
	\]
	所用线积分对几乎处处的 \(t\) 有限；这已足以给全族长度下确界一个有限上界。取 \(\rho=1\) 时，每条连接曲线的长度至少为 \(s\) 的总变化 \(L\)，水平弧达到它，面积为 \(2\pi L\)。所以第一式成立。

	对核心闭路族，固定 \(s\) 的圆圈都在族中，于是
	\[
		L L_\rho(\Gamma_{\mathrm{core}})^2
		\le2\pi\int_0^L\int_0^{2\pi}\rho(s,t)^2\,\mathrm{d}t\,\mathrm{d}s
		=2\pi A_\rho.
	\]
	反过来，在常密度 \(1\) 下，一条绕行次数为 \(n\ne0\) 的闭曲线提升到条带后，两端的 \(t\) 相差 \(2\pi n\)，所以其长度至少为 \(2\pi|n|\ge2\pi\)。固定 \(s\) 的一次圆圈实现 \(2\pi\)，得到第二式。平坦长度 \(|\mathrm{d}w|\) 拉回原环域即为 \(|\mathrm{d}z|/|z|\)，所以所述密度同时实现两个极值。
\end{proof}

由共形不变性，两个有限圆环域双全纯同构时模量相等，因为双全纯映射诱导基本群同构，将非零绕行族一一对应。反过来，模量相等就是半径比相等，缩放 \(z\mapsto(a'/a)z\) 已给出同构。因此模量完整分类这些圆环域。更一般地，已给定一个与圆环域共形同构的抽象环域时，其模量可由任一这种坐标计算，结果与坐标无关。

\begin{example}
	\label{g4:exm:0505:punctured-modulus}
	在穿孔圆盘 \(0<|z|<1\) 中，核心闭路族的极值长度为 \(0\)，从穿孔端到外边界的连接族极值长度为 \(\infty\)。对前者，任意 \(\varepsilon<|z|<1\) 中的核心族都是整个族的子族，所以
	\[
		0\le\operatorname{EL}(\Gamma_{\mathrm{core}})
		\le\frac{2\pi}{\log(1/\varepsilon)}\longrightarrow0.
	\]
	对后者，取密度 \(\rho_\varepsilon(z)=|z|^{-1}\) 于 \(\varepsilon<|z|<1\)，其余处取零。每条连接曲线都穿过该有限环带，长度至少为 \(\log(1/\varepsilon)\)，而面积为 \(2\pi\log(1/\varepsilon)\)。相应商趋于无穷大。这使穿孔圆盘自然具有无穷模量，并与任一有限模环域区别开来。
\end{example}

\begin{proposition}
	\label{g4:prop:0505:modulus-core-length}
	设 \(A(a,b)\) 为有限圆环域，\(M\) 为其模量，\(\ell\) 为其 Poincaré 度量中核心闭测地线的长度。则
	\[
		\ell=\frac\pi M,\qquad
		\operatorname{EL}(\Gamma_{\mathrm{core}})=\frac\ell\pi.
	\]
\end{proposition}
\begin{proof}
	\cref{g4:cor:0504:annulus-metric}给出 \(\ell=2\pi^2/\log(b/a)\)，而 \(M=\log(b/a)/(2\pi)\)。代入并结合式\eqref{g4:eq:0505:annulus-extremal-length}即得。
\end{proof}

这里的双曲度量是环域自身的完整 Poincaré 度量。如果环域嵌入另一个双曲曲面，后者度量在环域上的限制一般更小，由\cref{g4:thm:0504:poincare-complete}的收缩不等式得到这一比较。因此上式中的 \(\ell\) 不能未经比较就换成环境曲面中的某条测地线长度。

\subsection{退化曲面的长颈模型}
\label{g4:subsec:0505:04}

一个环域的模量很大，表示其连接两端的方向相对于绕行方向很长。先将对数圆柱按周长归一化为 \(1\)，便得到
\[
	C_M=(0,M)\times(\mathbb R/\mathbb Z).
\]
它与 \(\{e^{-2\pi M}<|z|<1\}\) 共形同构。平坦密度为 \(1\) 时，连接方向的最短长度为 \(M\)，核心长度为 \(1\)；若将面积归一化为 \(1\)，两者分别为 \(\sqrt M\) 与 \(1/\sqrt M\)。这些是同一共形结构在不同面积归一化下的长度，极值长度始终为 \(M\) 与 \(1/M\)。

\ProseParagraphBreak

也可以从全纯粘合看见同一个参数。固定 \(0<c<1\)，给定 \(0<|t|<c^2\)，将两张坐标圆盘中的环带
\[
	|t|/c<|z|<c,\qquad |t|/c<|w|<c
\]
按 \(zw=t\) 识别。换坐标 \(w=t/z\) 全纯且导数非零，因此它们形成一段环形颈部，模量为
\begin{equation}
	\label{g4:eq:0505:plumbing-modulus}
	M(t)=\frac1{2\pi}\log\frac{c^2}{|t|}.
\end{equation}
当 \(t\to0\) 时，连接族的极值长度趋于无穷大，核心族的极值长度趋于零。参数的辐角改变两端的相对角度；其模决定这里算出的颈部长宽比。形式上的中心方程 \(zw=0\) 由两张相交坐标盘组成，提示结点处的两条分支；对每个 \(t\ne0\)，上述构造和模量计算均发生在光滑环带上。

\begin{proposition}
	\label{g4:prop:0505:truncated-cylinder}
	设 \(\ell,R>0\)，考虑双曲圆柱中的对称开环带
	\[
		N_{\ell,R}=\{(r,t):-R<r<R,\ t\in\mathbb R/\ell\mathbb Z\},
		\qquad g=\mathrm{d}r^2+\cosh^2r\,\mathrm{d}t^2.
	\]
	其共形模量为
	\begin{equation}
		\label{g4:eq:0505:truncated-modulus}
		\operatorname{Mod}N_{\ell,R}
		=\frac2\ell\int_0^R\frac{\mathrm{d}u}{\cosh u}
		=\frac2\ell\arctan(\sinh R).
	\end{equation}
	固定 \(c>0\)，取 \(0<\ell<c\) 及 \(R_\ell=\operatorname{arcosh}(c/\ell)\)，使两条截口的长度均为 \(c\)。则
	\[
		\operatorname{Mod}N_{\ell,R_\ell}
		=\frac2\ell\arccos(\ell/c)
		=\frac\pi\ell-\frac2c+O(\ell^2)
		\qquad(\ell\downarrow0).
	\]
	以任一截口为起点的每个固定宽度区域，其度量光滑趋于一个尖点端的相应区域。
\end{proposition}
\begin{proof}
	令
	\[
		v(r)=\int_0^r\operatorname{sech}u\,\mathrm{d}u
		=\arctan(\sinh r).
	\]
	则 \(g=\cosh^2r\,(\mathrm{d}v^2+\mathrm{d}t^2)\)，所以这个环带共形于周长 \(\ell\)、高度 \(2v(R)\) 的平坦圆柱。由环域模量公式，其模量就是高度除以周长，得到式\eqref{g4:eq:0505:truncated-modulus}。特别地 \(R\to\infty\) 时，模量趋于 \(\pi/\ell\)，与整个完备圆柱的计算一致。

	截圈的长度为 \(\ell\cosh R\)，故所取 \(R_\ell\) 确实对应固定截口长 \(c\)。因为
	\[
		\sinh R_\ell=\sqrt{c^2/\ell^2-1},
		\qquad \arctan\sqrt{c^2/\ell^2-1}=\arccos(\ell/c),
	\]
	精确公式成立；在零点展开 \(\arccos x=\pi/2-x+O(x^3)\)，便得渐近式。

	在正半边从截口向核心取距离 \(s=R_\ell-r\)，并令 \(x=t/\ell\in\mathbb R/\mathbb Z\)。度量成为
	\[
		g=\mathrm{d}s^2+f_\ell(s)^2\,\mathrm{d}x^2,
		\qquad
		f_\ell(s)=\ell\cosh(R_\ell-s)
		=c\cosh s-\sqrt{c^2-\ell^2}\sinh s.
	\]
	对任意固定 \(S<\infty\)，当 \(\ell\) 足够小时有 \(R_\ell>S\)，且在 \(0\le s\le S\) 上 \(f_\ell\) 连同所有阶导数一致趋于 \(ce^{-s}\)。因此度量光滑趋于
	\[
		\mathrm{d}s^2+c^2e^{-2s}\,\mathrm{d}x^2,
	\]
	这正是\cref{g4:exm:0504:cusp}中 \(Y=1/c\) 的尖点端。负半边由 \(r\mapsto-r\) 对称得到。与此同时 \(R_\ell\to\infty\)，所以从任一截口到核心的距离趋于无穷。
\end{proof}

这一模型同时解释两个观察：缩短中间的闭测地线，会使环域模量发散；从两侧固定截口看进去，颈部越来越像两个无限长的尖点。若一个曲面上实际给定了这样的嵌入环带，模量与长度的上述计算就适用于这段环带。证明任意短测地线都拥有某个统一宽度的嵌入领圈，还需要控制不同提升之间的相交；那是全局双曲几何的另一项任务。
% 纲要条目排为编号小节标题，并进入目录与书签；说明段仍保留于纲要注释中。


\begin{chaptersummary}
	调和函数可以由圆盘内的 Poisson 积分局部改进；对所有满足边界条件的次调和候选取上包络，便得到 Dirichlet 解，边界障碍保证它取到指定值。挖去小盘再取极限产生 Green 函数。其对数奇点给共轭微分一个整数周期，负指数因此成为单值全纯函数；边界模长与简单零点进一步使这个函数成为双全纯圆盘映照。

	\ProseParagraphBreak

	一般统一化把这个有边界构造与拓扑穷竭接起来。单连通性使边界圆周可填为盘，非紧性保证可选择包含任意紧集的盘。固定基点导数后，单叶函数的增长估计控制正映射，四分之一圆盘及接近恒等的估计控制逆映射。极限既非退化又满射，得到平面与圆盘两种非紧模型；紧情形由亏格零及 Riemann--Roch 得到球面。万有覆盖将三种单连通模型转为一般曲面的群商。

	\ProseParagraphBreak

	圆盘度量沿覆盖下降，给双曲型曲面一个完备的曲率 \(-1\) 度量，全纯映射相对于它收缩；局部取等号对应全纯覆盖。极值长度则把曲线长度与面积一起比较，留下共形不变量。环域连接族与核心族的极值长度互为倒数，模量增大时核心的双曲长度缩短。显式圆柱和粘合模型使这类退化同时表现为长颈、短核心与两侧尖点。
\end{chaptersummary}

\BookExercises

\begin{exercise}
	\label{g4:ex:0501:blaschke-green}
	设
	\[
		B(z)=e^{\mathrm{i}\theta}\prod_{j=1}^N\frac{z-a_j}{1-\overline{a_j}z},
		\qquad a_j\in\mathbb D,\quad N\ge1,
	\]
	其中允许 \(a_j\) 重复。对任意 \(a\in\mathbb D\)，记 \(B^{-1}(a)=\{p_1,\ldots,p_k\}\)，并以 \(m_i\) 表示 \(B-a\) 在 \(p_i\) 的零点阶数。证明 \(\sum_i m_i=N\)，并证明
	\[
		G_{\mathbb D}(B(z),a)=\sum_{i=1}^k m_iG_{\mathbb D}(z,p_i).
	\]
\end{exercise}
\begin{solution}
	每个分母在闭圆盘上不为零，且 \(|B|=1\) 于单位圆周。对 \(|a|<1\)，Rouché 定理说明 \(B-a\) 与 \(B\) 在圆盘内有相同的零点总数 \(N\)，得到重数公式。

	左边在各 \(p_i\) 之外调和。在 \(p_i\) 附近写 \(B(z)-a=(z-p_i)^{m_i}c_i(z)\)，其中 \(c_i\) 全纯且不为零。由圆盘 Green 函数的显式表达，左边等于 \(-m_i\log|z-p_i|\) 加一条在 \(p_i\) 附近调和的函数。因此两边之差在所有 \(p_i\) 处都可去，成为整个圆盘上的调和函数。两边都连续趋于零边值：右边逐项如此；左边由 \(B\) 的边界模长及 \(a\) 严格位于圆盘内得到。最大值原则使差恒为零。这也说明拉回 Green 函数时，分歧以局部零点重数进入公式。
\end{solution}

\begin{exercise}
	\label{g4:ex:0502:robin-radius}
	设 \(U\) 为光滑边界闭拓扑盘的内部，\(p\in U\)，\(z(p)=0\) 为局部全纯坐标。在 \(p\) 附近写
	\[
		G_U(q,p)=-\log|z(q)|+H_z(q).
	\]
	设 \(r_z(U,p)\) 为该坐标下的共形半径。证明 \(H_z(p)=\log r_z(U,p)\)，并求换坐标 \(w=cz+O(z^2)\)、\(c\ne0\) 时两者的变化。若 \(\lambda_U(z)|\mathrm{d}z|\) 是 \(U\) 的 Poincaré 长度元，证明 \(\lambda_U(0)=2/r_z(U,p)\)。
\end{exercise}
\begin{solution}
	取归一化圆盘映照 \(f_U\)，使 \(f_U(p)=0\)、\(f_U'(p)>0\)。\cref{g4:thm:0502:bordered-disk}给出 \(G_U=-\log|f_U|\)。由于
	\[
		f_U(q)=f_U'(p)z(q)+O(z(q)^2),
	\]
	消去对数奇点后有 \(H_z(p)=-\log f_U'(p)\)。共形半径定义为 \(1/f_U'(p)\)，故得到所求等式。

	在新坐标下，
	\[
		H_w=H_z+\log|w/z|,
	\]
	而 \(w/z\) 在中心取值 \(c\)，所以 \(H_w(p)=H_z(p)+\log|c|\)，相应 \(r_w=|c|r_z\)。最后，Poincaré 度量由圆盘度量拉回，有
	\[
		\lambda_U(z)=\frac{2|f_U'(z)|}{1-|f_U(z)|^2}.
	\]
	在 \(z=0\) 代入即得 \(\lambda_U(0)=2f_U'(p)=2/r_z\)。换坐标时长度密度除以 \(|c|\)，与共形半径乘以 \(|c|\) 相容。
\end{solution}

\begin{exercise}
	\label{g4:ex:0503:radius-distance}
	设 \(\Omega\subsetneq\mathbb C\) 为非空单连通区域，\(p\in\Omega\)，记 \(\delta=\operatorname{dist}(p,\partial\Omega)\)。用通常坐标 \(z-p\) 定义共形半径 \(r(\Omega,p)\)。证明
	\[
		\delta\le r(\Omega,p)\le4\delta.
	\]
	分别给出两端取等号的例子，并证明 \(k(z)=z/(1-z)^2\) 把 \(\mathbb D\) 双全纯映到 \(\mathbb C\setminus(-\infty,-1/4]\)。
\end{exercise}
\begin{solution}
	统一化给出 \(\Omega\) 与圆盘或平面同构；后一种情形会给出整个平面到 \(\Omega\) 的单叶整函数。\cref{g4:thm:0503:noncompact-uniformization}证明中已由二次系数界推出单叶整函数只能仿射，其像应为整个平面，矛盾。因此取 \(f:\Omega\to\mathbb D\)，满足 \(f(p)=0\)、\(f'(p)>0\)。

	圆盘 \(D(p,\delta)\subset\Omega\)，对 \(\zeta\mapsto f(p+\delta\zeta)\) 用 Schwarz 引理，得 \(\delta f'(p)\le1\)，所以 \(\delta\le r\)。另一方面，令 \(\psi=f^{-1}\)，则 \((\psi(\zeta)-p)/r\) 单叶且导数为 \(1\)。\cref{g4:lem:0503:distortion}的四分之一圆盘结论给出 \(D(p,r/4)\subset\Omega\)，故 \(r\le4\delta\)。

	圆盘 \(D(p,R)\) 满足 \(r=\delta=R\)，给出左端等号。为处理右端，设 \(w=(1+z)/(1-z)\)，它把 \(\mathbb D\) 双全纯映到右半平面，且
	\[
		k(z)=\frac{w^2-1}{4}.
	\]
	平方把右半平面双全纯映到 \(\mathbb C\setminus(-\infty,0]\)：每个像点恰有一个实部为正的平方根。因此得到题中所述裂缝域。又 \(k(0)=0\)、\(k'(0)=1\)，该域在 \(0\) 的共形半径为 \(1\)，边界距离为 \(1/4\)，右端取等号。
\end{solution}

\begin{exercise}
	\label{g4:ex:0503:nonnested-exhaustion}
	设 \(X\) 为单连通非紧黎曼曲面，固定 \(p\) 及中心坐标。设 \(V_n\Subset X\) 均为光滑边界闭拓扑盘的内部，包含 \(p\)，并满足：对每个紧集 \(K\subset X\)，充分大的 \(n\) 均有 \(K\subset V_n\)。不要求 \(V_n\) 递增。令 \(F_n:V_n\to D_{r_n}\) 满足 \(F_n(p)=0\)、\(F_n'(p)=1\)。证明 \(r_n\to R\)，且 \(F_n\) 在紧集上收敛于正文的归一化统一化映照 \(F:X\to D_R\)。并说明单位圆盘归一化映照 \(f_n=F_n/r_n\) 在 \(R=\infty\) 时的极限。
\end{exercise}
\begin{solution}
	固定正文中的递增圆盘穷竭 \(U_m\)，其半径记为 \(s_m\uparrow R\)。对每个 \(m\)，充分大的 \(n\) 有 \(\overline U_m\subset V_n\)，所以半径单调性给出 \(r_n\ge s_m\)。若 \(R<\infty\)，则 \(F|_{V_n}\circ F_n^{-1}:D_{r_n}\to D_R\) 固定原点且导数为 \(1\)，Schwarz 引理经缩放给出 \(r_n\le R\)。于是 \(r_n\to R\)。若 \(R=\infty\)，对任意有限下界选择充分大的 \(s_m\)，同样得到 \(r_n\to\infty\)。

	对固定 \(m\)，以 \(\phi_m:\mathbb D\to U_m\) 表示满足 \(\phi_m(0)=p\)、\(\phi_m'(0)=s_m\) 的映照。则充分大的 \(n\) 使 \(F_n\circ\phi_m/s_m\) 在圆盘上单叶且归一化。增长估计给出 \(U_m\) 的任一紧集上的统一界。因此任意子列都可抽出在 \(X\) 的紧集上收敛的子列，其极限 \(G\) 满足 \(G(p)=0\)、\(G'(p)=1\)，并且单叶。确实，若两个不同点有相同的极限值，在两点的不交小坐标盘中取边界上不出现该值的圆周，Rouché 定理使充分大的映照在两个盘中都取到该值，与单射性矛盾。

	若 \(R<\infty\)，\(|F_n|<r_n\le R\) 使 \(G(X)\subset D_R\)。复合 \(G\circ F^{-1}:D_R\to D_R\) 固定原点且导数为 \(1\)，故由 Schwarz 等号条件恒等。若 \(R=\infty\)，该复合是单叶整函数，故仿射；归一化同样迫使它恒等。因此每个收敛子列的极限都是 \(F\)，正规性使整列收敛。抛物情形中，\(F_n\) 在每个紧集上有界而 \(r_n\to\infty\)，所以 \(f_n\to0\) 局部一致；只把基点送到原点确实可能留下常值极限。
\end{solution}

\begin{exercise}
	\label{g4:ex:0503:two-omitted-values}
	证明 \(\mathbb C\setminus\{0,1\}\) 为双曲型曲面。进一步证明：如果整函数 \(f:\mathbb C\to\mathbb C\) 的像遗漏两个不同的复数，则 \(f\) 为常值。论证中说明如何提升映射，以及提升为什么仍全纯。
\end{exercise}
\begin{solution}
	第一卷第5.2节命题5.2.10计算了平面去掉有限个点的基本群：每个缺失点的小圆周给出一个自由生成元。因此 \(X=\mathbb C\setminus\{0,1\}\) 的基本群为二生成元自由群 \(F_2\)，特别地不交换。若 \(X\) 的万有覆盖是球面，则覆盖投影的像紧，与 \(X\) 非紧矛盾。若万有覆盖是平面，\secref{g4:sec:0503}的商分类说明覆盖变换都是平移，\(X\) 为 \(\mathbb C\)、\(\mathbb C^\times\) 或复环面之一；它们的基本群分别为 \(0\)、\(\mathbb Z\)、\(\mathbb Z^2\)，都交换。故只剩圆盘覆盖，\(X\) 为双曲型。

	若 \(f\) 遗漏 \(a\ne b\)，则 \(F=(f-a)/(b-a)\) 是从 \(\mathbb C\) 到 \(X\) 的全纯映射。取全纯万有覆盖 \(\pi:\mathbb D\to X\)，并指定 \(F(0)\) 的一个提升。第一卷第5.3节定理5.3.10的提升判据要求 \(F_*\pi_1(\mathbb C,0)\subseteq\pi_*\pi_1(\mathbb D)\)；两边均为零，所以存在连续提升 \(\widetilde F:\mathbb C\to\mathbb D\)。在每个足够小的源邻域中，\(\widetilde F\) 等于 \(F\) 后接 \(\pi\) 的某个局部全纯逆，因此提升全纯。它是有界整函数，由 Liouville 定理为常数。于是 \(F=\pi\widetilde F\) 及 \(f\) 都为常数。
\end{solution}

\begin{exercise}
	\label{g4:ex:0504:power-coverings}
	设 \(n\ge2\)。直接用密度公式证明 \(P_n(z)=z^n\) 从穿孔单位圆盘到自身是 Poincaré 度量的局部等距。对有限环域 \(A(a,b)\)，证明它给出到 \(A(a^n,b^n)\) 的 \(n\) 重全纯覆盖，并比较两者的模量与核心闭测地线长度。解释局部等距与目标核心长度改变为何相容。
\end{exercise}
\begin{solution}
	在穿孔圆盘上，长度密度为 \(\lambda(z)=1/(|z|\log(1/|z|))\)。直接代入得
	\[
		\lambda(z^n)n|z|^{n-1}
		=\frac{n|z|^{n-1}}{|z|^n n\log(1/|z|)}
		=\lambda(z).
	\]
	每个非零目标点有 \(n\) 个不同的 \(n\) 次方根，局部可选全纯根支，因此映射为 \(n\) 重覆盖。

	有限环域中的同一论证给出覆盖，且
	\[
		M'=\frac1{2\pi}\log\frac{b^n}{a^n}=nM.
	\]
	由\cref{g4:prop:0505:modulus-core-length}，相应核心长度满足 \(\ell'=\pi/M'=\ell/n\)。也可直接代入式\eqref{g4:eq:0504:annulus-density}：新对数宽度为 \(n\log(b/a)\)，正弦中的比值不变，乘 \(n|z|^{n-1}\) 后恢复源密度。

	源核心 \(|z|=\sqrt{ab}\) 映到目标核心 \(|w|=(ab)^{n/2}\)，但绕行 \(n\) 次。按参数曲线计算，其像的总长度为 \(n\ell'=\ell\)，恰等于源长度。局部等距保持的是带绕行次数的曲线长度，目标核心作为一次简单闭曲线的长度则为 \(\ell/n\)。
\end{solution}

\begin{exercise}
	\label{g4:ex:0504:branched-density}
	仍取 \(P_n(z)=z^n\)，现在把它看成整个单位圆盘到自身的映射。计算拉回 Poincaré 度量的长度密度 \(\mu_n\)，证明它在 \(0\) 为零，而在每个内点严格小于圆盘密度。证明两密度之比沿半径趋向边界时趋于 \(1\)。再证明 \(\mu_n^2|\mathrm{d}z|^2\) 在穿孔圆盘上曲率为 \(-1\)，却不完备。
\end{exercise}
\begin{solution}
	拉回密度为
	\[
		\mu_n(z)=\frac{2n|z|^{n-1}}{1-|z|^{2n}},
	\]
	所以 \(\mu_n(0)=0\)。当 \(0<r=|z|<1\) 时，
	\[
		\frac{\mu_n(z)}{\lambda_{\mathbb D}(z)}
		=\frac{nr^{n-1}}{1+r^2+\cdots+r^{2n-2}}<1.
	\]
	严格不等式来自算术--几何平均不等式：分母的 \(n\) 项几何平均为 \(r^{n-1}\)，且 \(r<1\)、\(n\ge2\) 时这些项不全相等。在 \(0\) 处比值为零；令 \(r\uparrow1\)，分子与分母均趋于 \(n\)，所以比值趋于 \(1\)。边界上的极限等号并不是 Schwarz--Pick 在内点的等号。

	在 \(z\ne0\) 处，\(P_n'\ne0\)，拉回为局部等距，故曲率为 \(-1\)。但是沿径向线段从某个 \(r_0>0\) 走向穿孔端，长度为
	\[
		\int_0^{r_0}\frac{2nr^{n-1}}{1-r^{2n}}\,\mathrm{d}r
		=\int_0^{r_0^n}\frac2{1-s^2}\,\mathrm{d}s<\infty.
	\]
	取径向趋于零的点列，尾部长度趋零使它成为 Cauchy 列，却在穿孔圆盘中没有极限。因此这个度量不完备；它也不是穿孔圆盘自身的 Poincaré 度量，后者到穿孔端的距离无限。
\end{solution}

\begin{exercise}
	\label{g4:ex:0501:irregular-boundary}
	设 \(\Omega=\{0<|z|<1\}\)，在其拓扑边界 \(\{0\}\cup\{|z|=1\}\) 上给定连续函数 \(f(0)=1\)、\(f|_{|z|=1}=0\)。证明不存在在 \(\Omega\) 调和且在 \(\overline\Omega\) 连续取该边值的函数。若仍对所有连续至闭包、内部次调和且边值不超过 \(f\) 的函数取上包络，求这个包络，并指出其在边界的行为。此题合并并展开了\parencite[第6章第4.2节，第249页及第251页习题]{Ahlfors1979Complex}中的反例与包络问题。
\end{exercise}
\begin{solution}
	若调和解存在，闭包紧性使它有界。\cref{g4:thm:0501:green}证明中给出的有界调和函数可去奇点论证，使它在 \(0\) 处延拓为调和函数。延拓后的函数在单位圆盘边界为零，由最大值原则恒为零，与 \(u(0)=1\) 矛盾。

	设 \(v\) 是题中的一个次调和候选，它在闭包上有上界 \(M\)。固定 \(\varepsilon>0\)，取 \(\delta\) 足够小使 \(\varepsilon\log(1/\delta)\ge\max(M,0)\)。在 \(\delta<|z|<1\) 中，将 \(v\) 与调和函数 \(\varepsilon\log(1/|z|)\) 比较：两条边界上前者均不超过后者，故内部也如此。对任一固定 \(z\in\Omega\) 可任意缩小 \(\delta<|z|\)，于是 \(v(z)\le\varepsilon\log(1/|z|)\)。再令 \(\varepsilon\downarrow0\)，得到 \(v\le0\)。常数零本身是候选，所以包络在 \(\Omega\) 内恰为零。它趋向穿孔端时仍为零，没有取得指定的边值 \(1\)。边界障碍正是 Perron 构造中用于排除这种现象的条件。
\end{solution}

\begin{exercise}
	\label{g4:ex:0505:torus-extremal-length}
	设 \(E_\tau=\mathbb C/(\mathbb Z+\tau\mathbb Z)\)，\(Y=\operatorname{Im}\tau>0\)，且整数 \(p,q\) 互素。记 \(v=p+q\tau\)，令 \(\Gamma_v\) 为提升到平面后两端之差为 \(v\) 的全部有向闭曲线。证明
	\[
		\operatorname{EL}_{E_\tau}(\Gamma_v)=\frac{|v|^2}{Y}.
	\]
	若一个有限模环域共形嵌入 \(E_\tau\)，其一次核心曲线代表 \(v\) 或 \(-v\)，证明该环域的模量不超过 \(Y/|v|^2\)，并构造达到此界的嵌入环域。
\end{exercise}
\begin{solution}
	由 \(p,q\) 互素，可取整数 \(r,s\) 使 \(ps-qr=1\)。令 \(w=r+s\tau\)，则 \(v,w\) 是原格的一组正向整数基，且 \(\operatorname{Im}(\overline v w)=Y\)。以 \(z=v\sigma+wt\)、\(0\le\sigma,t\le1\) 表示基本平行四边形，面积元为 \(Y\,\mathrm{d}\sigma\,\mathrm{d}t\)。固定 \(t\) 的线段在商上是 \(\Gamma_v\) 的成员。对任一可容许密度 \(\rho\)，Cauchy--Schwarz 与 Fubini 给出
	\[
		L_\rho(\Gamma_v)^2
		\le |v|^2\int_0^1\rho(v\sigma+wt)^2\,\mathrm{d}\sigma
	\]
	对几乎处处的 \(t\) 成立，积分后得到 \(L_\rho^2\le |v|^2A_\rho/Y\)。常密度 \(1\) 的面积为 \(Y\)，而任一所述曲线的平面提升两端相差 \(v\)，长度至少为 \(|v|\)；直线闭路实现此下界。因此上界恰可达到。

	嵌入环域中取与 \(v\) 同向的一次核心闭路族，它是 \(\Gamma_v\) 的子族。环域计算同样适用于一次绕行族，其极值长度为模量 \(M\) 的倒数。由曲线族单调性及开子域补零不改变极值长度，
	\[
		\frac{|v|^2}{Y}=\operatorname{EL}(\Gamma_v)\le\frac1M,
	\]
	所以 \(M\le Y/|v|^2\)。

	在坐标 \(\zeta=z/v\) 中，格为 \(\mathbb Z+(w/v)\mathbb Z\)，且 \(H=\operatorname{Im}(w/v)=Y/|v|^2\)。从环面删去实轴的像，剩余区域由条带 \(0<\operatorname{Im}\zeta<H\) 模整数平移表示；条带内部没有额外格识别，因为 \(w/v\) 的平移改变虚部 \(H\)。指数 \(e^{2\pi\mathrm{i}\zeta}\) 将其识别为 \(A(e^{-2\pi H},1)\)，模量为 \(H\)，达到所求界。
\end{solution}

\begin{exercise}
	\label{g4:ex:0505:annulus-map-degree}
	设 \(A_j=A(a_j,b_j)\)，\(0<a_j<b_j<\infty\)，且 \(f:A_1\to A_2\) 全纯。设 \(f\) 在核心闭路上诱导的绕数为整数 \(n\ne0\)。证明
	\[
		|n|\operatorname{Mod}A_1\le\operatorname{Mod}A_2.
	\]
	反过来，若这条不等式成立，构造具有指定绕数 \(n\) 的全纯映射。特别说明等号时如何选择一个全纯覆盖。
\end{exercise}
\begin{solution}
	在源环域的对数覆盖上写 \(z=a_1e^{s+\mathrm{i}t}\)，其中 \(0<s<L_1=\log(b_1/a_1)\)，\(t\in\mathbb R\)。由于条带单连通而 \(f\ne0\)，可选 \(f(a_1e^{s+\mathrm{i}t})\) 的全纯对数 \(u(s,t)+\mathrm{i}v(s,t)\)。其实部 \(u=\log|f|\) 以 \(2\pi\) 为 \(t\) 周期，而绕数条件给出
	\[
		v(s,t+2\pi)-v(s,t)=2\pi n.
	\]
	Cauchy--Riemann 方程为 \(u_s=v_t\)，故平均值
	\[
		\overline u(s)=\frac1{2\pi}\int_0^{2\pi}u(s,t)\,\mathrm{d}t
	\]
	满足 \(\overline u'(s)=n\)。因此 \(\overline u(s)=ns+c\)。它始终介于 \(\log a_2\) 与 \(\log b_2\) 之间，令 \(s\) 分别趋于两端，得到 \(|n|L_1\le\log(b_2/a_2)=L_2\)。除以 \(2\pi\) 即得所求。

	若 \(n>0\)，取 \(f(z)=c_0(z/a_1)^n\)，其中 \(c_0>0\)。它的对数模取值区间为 \((\log c_0,\log c_0+nL_1)\)。由 \(nL_1\le L_2\)，可以将此区间放在 \((\log a_2,\log b_2)\) 中。若 \(n<0\)，取同一公式 \(f(z)=c_0(z/a_1)^n\)，区间方向反转，选择 \(\log c_0\) 靠近上端即可。两种情形的绕数均为 \(n\)。等号时令两个区间端点分别相等，映射即为 \(|n|\) 重覆盖；每个目标点的根都位于源环域，且导数不为零。
\end{solution}

\begin{exercise}
	\label{g4:ex:0501:period-energy}
	设 \(A=A(a,b)\)、\(M=\operatorname{Mod}A\)，并固定 \(c\in\mathbb R\)。在所有光滑闭实一形式 \(\alpha\) 中，要求沿正向核心圆周的周期为 \(c\)，且 \(\int_A|\alpha|^2\,\mathrm{d}A<\infty\)。证明其能量下确界为 \(c^2M\)，并求唯一达到下确界的形式。这里能量相对于任一共形度量计算。
\end{exercise}
\begin{solution}
	一形式能量在共形换度量下不变，因为余切范数乘逆平方因子，面积乘平方因子。用对数坐标 \(z=a e^{s+\mathrm{i}t}\)，令 \(0<s<L=\log(b/a)\)、\(t\in\mathbb R/2\pi\mathbb Z\)，并写
	\[
		\alpha=A(s,t)\,\mathrm{d}s+B(s,t)\,\mathrm{d}t.
	\]
	闭性为 \(B_s=A_t\)。对 \(t\) 积分，周期 \(\int_0^{2\pi}B(s,t)\,\mathrm{d}t\) 与 \(s\) 无关，所以处处等于 \(c\)。Cauchy--Schwarz 给出
	\[
		\int_0^{2\pi}B(s,t)^2\,\mathrm{d}t\ge\frac{c^2}{2\pi}.
	\]
	再对 \(s\) 积分，并加上非负的 \(A^2\) 项，得到
	\[
		\int_A|\alpha|^2\,\mathrm{d}A
		=\int_0^L\int_0^{2\pi}(A^2+B^2)\,\mathrm{d}t\,\mathrm{d}s
		\ge\frac{c^2L}{2\pi}=c^2M.
	\]
	等号要求 \(A=0\)，且每条圆周上 \(B\) 为常数；固定周期使这个常数为 \(c/(2\pi)\)。由光滑性，几乎处处的等式扩展到每一点。因此唯一极小形式是
	\[
		\alpha=\frac c{2\pi}\,\mathrm{d}t
		=\frac c{2\pi}\,\mathrm{d}\arg z.
	\]
	它确实闭、具有指定周期和有限能量。当 \(c\ne0\) 时它不是恰当形式；周期约束保留了环域的整体信息。
\end{solution}